% concatenated from paper.tex
\documentclass{svjour3}


\usepackage{amsfonts}
\usepackage{amsmath}
\usepackage{graphicx}
\usepackage{geometry}
\usepackage{csquotes}
\usepackage{amssymb}
\usepackage{algorithm}
\usepackage{algorithmicx}
\usepackage{algpseudocode}

\algrenewcommand\algorithmicrequire{\textbf{Input:}}
\algrenewcommand\algorithmicensure{\textbf{Output:}}



\usepackage{color}
\usepackage[hidelinks]{hyperref}
\usepackage{mathrsfs}
\setcounter{MaxMatrixCols}{10}


% \geometry{hmargin=3cm,vmargin=3cm}
% \newtheorem{theorem}{Theorem}[section]
% \newtheorem{acknowledgement}[theorem]{Acknowledgement}
% \newtheorem{axiom}[theorem]{Axiom}
% \newtheorem{case}[theorem]{Case}
% \newtheorem{claim}[theorem]{Claim}
% \newtheorem{conclusion}[theorem]{Conclusion}
% \newtheorem{condition}[theorem]{Condition}
% \newtheorem{conjecture}[theorem]{Conjecture}
% \newtheorem{corollary}[theorem]{Corollary}
% \newtheorem{criterion}[theorem]{Criterion}
% \newtheorem{definition}[theorem]{Definition}
% \newtheorem{example}[theorem]{Example}
% \newtheorem{exercise}[theorem]{Exercise}
% \newtheorem{lemma}[theorem]{Lemma}
% \newtheorem{notation}[theorem]{Notation}
% \newtheorem{problem}[theorem]{Problem}
% \newtheorem{proposition}[theorem]{Proposition}
% \newtheorem{remark}[theorem]{Remark}
% \newtheorem{solution}[theorem]{Solution}
% \newtheorem{summary}[theorem]{Summary}
% \newenvironment{proof}[1][Proof]{\noindent\textbf{#1.} }{\ \rule{0.5em}{0.5em}}

\newcommand{\R}{\mathbb{R}} % Reals
\newcommand{\N}{\mathbb{N}} % Naturals
\newcommand{\de}{\mathrm{d}}
\newcommand{\T}{\mathcal{T}}

\renewcommand{\P}{\mathbb{P}} % Polynomials
\renewcommand{\de}{\mathrm{d}} % Differential
\newcommand\smallO{
  \mathchoice
    {{\scriptstyle\mathcal{O}}}% \displaystyle
    {{\scriptstyle\mathcal{O}}}% \textstyle
    {{\scriptscriptstyle\mathcal{O}}}% \scriptstyle
    {\scalebox{.7}{$\scriptscriptstyle\mathcal{O}$}}%\scriptscriptstyle
  }


\author{Ludovico Bruni Bruno         \and
        Francesco Dell'Accio \and 
        Wolfgang Erb \and 
        Federico Nudo
}

%\authorrunning{Short form of author list} % if too long for running head

\institute{Ludovico Bruni Bruno \at
              Department of Mathematics \enquote{Tullio Levi-Civita}, University of Padova, Italy \\
              Istituto Nazionale di Alta Matematica, Roma, Italy
              \email{ludovico.brunibruno@unipd.it}         
                %  \\
%             \emph{Present address:} of F. Author  %  if needed
           \and
            Francesco Dell'Accio \at
             Department of Mathematics and Computer Science, University of Calabria, Rende (CS), Italy\\ 
             Istituto per le Applicazioni del Calcolo 'Mauro Picone', Naples Branch, C.N.R. National Research Council of Italy, Napoli, Italy
             \email{francesco.dellaccio@unical.it} 
         \and 
             Wolfgang Erb \at
              Department of Mathematics \enquote{Tullio Levi-Civita}, University of Padova, Italy \\
              \email{wolfgang.erb@unipd.it}
         \and 
            Federico Nudo \at
              Department of Mathematics \enquote{Tullio Levi-Civita}, University of Padova, Italy \\
              \email{federico.nudo@unipd.it}
}

\begin{document}
% \bigskip

% \bigskip 
% \begin{frontmatter}

\title{Bivariate polynomial histopolation techniques on Padua, Fekete and Leja triangles}
%% Group authors per affiliation:

% \author[address-Pd,address-IN]{Ludovico Bruni Bruno\corref{corrauthor}}\cortext[corrauthor]{Corresponding author}
% \ead{ludovico.brunibruno@unipd.it}

% \author[address-CS]{Francesco Dell'Accio}
% \ead{francesco.dellaccio@unical.it}

% \author[address-Pd]{Wolfgang Erb}
% \ead{wolfgang.erb@unipd.it}


% \author[address-Pd]{Federico Nudo}%\corref{corrauthor}}\cortext[corrauthor]{Corresponding author}
% \ead{federico.nudo@unipd.it}

% \address[address-Pd]{Department of Mathematics \enquote{Tullio Levi-Civita}, University of Padova, Italy}

% \address[address-IN]{Istituto Nazionale di Alta Matematica \enquote{Francesco Severi}, Roma, Italy}

% \address[address-CS]{Department of Mathematics and Computer Science, University of Calabria, Rende (CS), Italy}
\date{}

\maketitle


\begin{abstract}
This paper explores the reconstruction of a real-valued function $f$ defined over a domain $\Omega \subset \mathbb{R}^2$ using bivariate polynomials that satisfy triangular histopolation conditions. More precisely, we assume that only the averages of $f$ over a given triangulation $\mathcal{T}_N$ of $\Omega$ are available and seek a bivariate polynomial that approximates $f$ using a histopolation approach, potentially flanked by an additional regression technique. This methodology relies on the selection of a subset of triangles $\mathcal{T}_M \subset \mathcal{T}_N$ for histopolation, ensuring both the solvability and the well-conditioning of the problem. The remaining triangles can potentially be used to enhance the accuracy of the polynomial approximation through a simultaneous regression. We will introduce histopolation and combined histopolation-regression methods using the Padua points, discrete Leja sequences, and approximate Fekete nodes. The proposed algorithms are implemented and evaluated through numerical experiments that demonstrate their effectiveness in function approximation.
\keywords{Polynomial histopolation \and Polynomial histopolation-regression \and Padua points \and Discrete Leja sequences \and Approximate Fekete points }
\subclass{33F05 \and 41A05 \and 41A10}
\end{abstract}

\section{Introduction}

Let $f$ be a real-valued function defined on a polygonal domain $\Omega \subset \R^2$. We assume that this domain is partitioned into a set of triangles $ \mathcal{T}_N=\left\{t_1, \ldots, t_N \right\} $ 
such that
\begin{equation*} \Omega = \bigcup_{i=1}^N t_i, \quad \text{ and } \quad   \left\lvert t_i \cap t_j \right\rvert = 0, \quad i\neq j. 
\end{equation*}
The main goal of this paper is to explore approximations of the function $f$ through a bivariate polynomial $p(x,y)$, assuming that only the averages
\begin{equation} \label{eq:mui} \mu_i(f) = \frac{1}{|t_i|}\int_{t_i} f(x,y) \de x \de y , \quad i \in \{ 1, \ldots, N\},
\end{equation}
of $f$ over the triangles $t_i$ are available. Such problems arise in various fields, including histosplines~\cite{Schoenberg}, optimal transport~\cite{OptimalTransport}, and mimetic methods~\cite{BeiraoMimetic}. 
The first obstacle in the resolution of such a problem emerges when the number $N$ of data values does not match the dimension $ M \doteq \operatorname{dim} \left(\P_m\left(\mathbb{R}^2\right)\right)$ of the chosen polynomial space, making exact histopolation impossible. %In this case, we shall consider a subsampling $ \mathcal{T}_M $ of $ \mathcal{T}_N $ and thus the question of correctly selecting such supports arises. %We propose a strategy involving sub-sampling and regression to ensure stability and accuracy in this framework.
%In fact, let  $\P_m\left(\mathbb{R}^2\right)$ be the space of bivariate polynomials of total degree $m$, whose dimension is 
%\begin{equation}\label{M}
%    M \doteq \operatorname{dim} \P_m\left(\mathbb{R}^2\right) = \binom{m+2}{2}. 
%\end{equation}
Even if $N = M$ is satisfied, the 
\emph{histopolation} problem which can be formulated as (cf. \cite{Robidoux})
\begin{equation} \label{eq:histopol} \text{Given a function $f$, find a polynomial $p$ such that $\mu_i(p) = \mu_i(f)$, $i \in \{ 1,\dots,M\}$,} \end{equation}
might not have a unique solution in the space $\P_m \left(\mathbb{R}^2\right)$ of bivariate polynomials of total degree less or equal to $m$. Furthermore, even if a unique solution exists, the histopolation problem \eqref{eq:histopol} may be ill-conditioned~\cite{RothThesis}. Such an ill-conditioning, which is typically due to a disadvantageous selection of triangles, causes inaccuracies in the polynomial approximation of the function $f$ and leads to Runge-like phenomena as, for instance, displayed in Figure~\ref{fig:Runge}. 
\begin{figure}[ht]
	\centering
	\includegraphics[width=0.49\textwidth]{Runge10.eps}
	\includegraphics[width=0.49\textwidth]{RandomRunge10.eps}
	\caption{Ill-conditioning of histopolation: the Runge function $ f(x,y) = \frac{1}{1+10(x^2 + y^2)} $ (left) and its polynomial reconstruction from average data on a random selection of triangles (right) in $[-1,1]^2$, using a total polynomial degree $ m = 5 $.}
	\label{fig:Runge}
\end{figure}

To overcome these issues, we refine the triangulation $\T_N$ and sub-sample $\mathcal{T}_N$ on a smaller set
$\T_M \subset \T_N$ of $M < N$ triangles, verifying that the histopolation problem is solvable, i.e., that the collection $ \T_M $ is unisolvent for the polynomial space $\P_m\left(\mathbb{R}^2\right)$, and well-conditioned. Furthermore, additional data values on the remaining triangles may be used to enhance the approximation accuracy through a simultaneous regression. 
This leads to the combined \emph{histopolation-regression} problem 
\begin{equation}\label{consleastsquares} \text{Find a polynomial $p$ such that } \begin{cases} \mu_i(p) = \mu_i(f), & i \in \{ 1, \ldots, M\}, \\ \mu_{M+j}(p) \approx \mu_{M+j}(f), & j \in \{ 1, \dots, N-M\}, \end{cases} \end{equation}
where $p$ lies in a larger polynomial space $\P_d\left(\mathbb{R}^2\right)$, $d > m$, and the set $\T_N$ is reordered so that its first elements are those of $\T_M$. 
%By setting the triangulation $\T_N$ and a suitable subset $\T_M\subset \T_N$, w
%We then define the relative \textit{histopolation-regression} operator as an operator
%\begin{equation*} \begin{array}{rcl} \Pi: f &\mapsto& \Pi[f], \end{array} \end{equation*}
%$$ \Pi: \ C (\Omega) \to \P_d\left(\mathbb{R}^2\right)  $$
%such that $\Pi(f)$ satisfies~\eqref{consleastsquares} for a suitable chosen polynomial space $\mathbb{P}_d\left(\mathbb{R}^2\right)$ . 
In this paper, we will explore three different strategies for histopolation and combined histopolation-regression. The first one is based on the \emph{Padua points}. These set of nodes are well-established in bivariate interpolation \cite{Bos;2006:BLI,bos2007bivariate} and exploit the \emph{Dubiner metric} \cite{Dubiner}
$$ \mathrm{dist} (\xi, \eta) \doteq \max \left\{ \left|\arccos\left(\xi_x\right) - \arccos\left(\eta_x\right) \right|, \left| \arccos\left(\xi_y\right) - \arccos\left(\eta_y\right) \right| \right\} $$
for points $ \xi = \left(\xi_x, \xi_y\right) $ and $ \eta = \left(\eta_x,\eta_y\right) $ in $ \Omega = [-1,1]^2 $. It is in fact conjectured in~\cite{CDMV} that equidistributed nodes with respect to such a distance are near-optimal for bivariate interpolation. The  relation between interpolation on the Padua points and histopolation on the Padua triangles is established through the mean value theorem and the results of Proposition \ref{prop:stability}. Indeed, while the mean value theorem allows to convert the histopolation problem into an interpolation problem which can be interpreted as a perturbation of the interpolation on the Padua points, Proposition \ref{prop:stability} ensures that the solution of the histopolation problem on the Padua triangles converges to the interpolant on the Padua points if the triangulation gets increasingly refined in a uniform sense. For consistency with the definition of the Padua points \cite{PaduaComputational}, we will generally consider histopolation problems on the domain $ \Omega = [-1,1]^2 $. Moreover, we will compare the histopolation results on the Padua triangles with two additional numerical triangle selection strategies that are inspired by established techniques for generating discrete Leja sequences and approximate Fekete points in classical polynomial interpolation. 

%To overcome this limitation, we propose two additional methods, inspired by techniques for generating discrete Leja and Fekete points in classical interpolation problems. 

\noindent \textbf{Outline of the paper}. The paper is organized as follows. In Section~\ref{sec2}, we extend the concept of polynomial interpolation on Padua points to bivariate polynomial histopolation on triangles. For this, we will present a histopolation procedure over a triangulation of the domain $[-1,1]^2$, together with an enhancement strategy based on an additional regression. We will establish bounds for the respective Lebesgue constants, i.e., the norms of the histopolation operator and the histopolation-regression operator. In Section~\ref{sec3new}, we will explore two additional histopolation-regression techniques for bivariate domains in which the numerical extraction of the triangles follows the strategy for the calculation of the approximate Fekete points and Leja sequences. These additional procedures will be used for the comparisons in Section~\ref{sec4}, where we will provide several numerical results that demonstrate the accuracy of the histopolation method based on the Padua triangles.






\section{Histopolation-regression techniques based on Padua triangles}\label{sec2}

\subsection{Determination of Padua triangles}
For the domain $ \Omega = [-1,1]^2 $, the set of Padua points $X_M^{\mathrm{Pad}}$ of size $ M = (m+1)(m+2)/2 $, $m \in \mathbb{N}$, is defined as
\begin{equation}\label{PaduaPoints}
	X_M^{\mathrm{Pad}}=\left\{\left((-1)^{i+j}\cos\left(\frac{\pi j}{m+1}\right),(-1)^{i+j}\cos\left(\frac{i\pi}{m}\right)\right)\, : \, 0\le i+j\le m\right\} .
\end{equation}
Given a triangulation $ \mathcal{T}_N = \left\{ t_i\right\}_{i=1}^N $ of size $N \geq M$, we define a corresponding set of Padua triangles as 
\begin{equation*}\label{PaduaTriangles}
	\mathcal{T}^{\mathrm{Pad}}_M \doteq \left\{ t \in \mathcal{T}_N \, :\,  \boldsymbol{x}^{\mathrm{Pad}} \in t \text{ for some }\boldsymbol{x}^{\mathrm{Pad}} \in 	X_M^{\mathrm{Pad}}  \right\}.
\end{equation*}


Since the set of nodes $ X_M^{\mathrm{Pad}} $ is unisolvent for the space of polynomials $ \P_m\left(\mathbb{R}^2\right)$, see~\cite{Bos;2006:BLI,bos2007bivariate}, we are interested in obtaining an attribution map  
\begin{equation} \label{eq:PPtoPT}
\varphi: \quad X^{\mathrm{Pad}}_M \ni \boldsymbol{x}^{\mathrm{Pad}} \mapsto t_i \in \mathcal{T}^{\mathrm{Pad}}_M
\end{equation}
which is not only well-defined but also injective, such that the number of Padua triangles in $\mathcal{T}^{\mathrm{Pad}}_M$ corresponds again to the dimension $M = (m+1)(m+2)/2$ of the space $\mathbb{P}_m\left(\mathbb{R}^2\right)$ of bivariate polynomials of total degree $m$.  For any point $ \boldsymbol{x} \in \R^2 $ and triangle $ t_i$ with vertices $\boldsymbol{v}^{(i)}_1, \boldsymbol{v}^{(i)}_2, \boldsymbol{v}^{(i)}_3,$ one has that $ \boldsymbol{x} \in t_i$ if and only if the barycentric coordinates of the triangle $t_i$ at the points $\boldsymbol{x}$ are nonnegative, that is
\begin{eqnarray*}
\lambda_{1}^{(i)}({\boldsymbol{
x}})&:=&\frac{\left(\boldsymbol{x}-\boldsymbol{v}^{(i)}_3\right)\times\left(\boldsymbol{v}^{(i)}_2-\boldsymbol{v}^{(i)}_3\right)}{\left(\boldsymbol{v}^{(i)}_1-\boldsymbol{v}^{(i)}_3\right)\times\left(\boldsymbol{v}^{(i)}_2-\boldsymbol{v}^{(i)}_3\right)}\ge0, 
\\ 
\lambda_{2}^{(i)}({\boldsymbol{
x}})&:=&\frac{\left(\boldsymbol{x}-\boldsymbol{v}^{(i)}_3\right)\times\left(\boldsymbol{v}^{(i)}_3-\boldsymbol{v}^{(i)}_1\right)}{\left(\boldsymbol{v}^{(i)}_1-\boldsymbol{v}^{(i)}_3\right)\times\left(\boldsymbol{v}^{(i)}_2-\boldsymbol{v}^{(i)}_3\right)}\ge0, \\
\lambda_{3}^{(i)}({\boldsymbol{
x}})&:=&\frac{\left(\boldsymbol{x}-\boldsymbol{v}^{(i)}_1\right)\times\left(\boldsymbol{v}^{(i)}_1-\boldsymbol{v}^{(i)}_2\right)}{\left(\boldsymbol{v}^{(i)}_1-\boldsymbol{v}^{(i)}_3\right)\times\left(\boldsymbol{v}^{(i)}_2-\boldsymbol{v}^{(i)}_3\right)}\ge0.
\end{eqnarray*}
Checking these conditions is computationally expensive. To be more cost-efficient, we can apply a \emph{point-in-triangle} algorithm \cite{HormannAgathosPIP} for determining whether a Padua point $\boldsymbol{x}^{\mathrm{Pad}}$ belongs to some $ t_i \in \mathcal{T}_N$. For this, consider a point $ \boldsymbol{x} \in \Omega $ and a triangle $ t_i $ of vertices $ \boldsymbol{v}^{(i)}_1, \boldsymbol{v}^{(i)}_2, \boldsymbol{v}^{(i)}_3 $. Then $ \boldsymbol{x} $ defines three (possibly degenerate) triangles by replacing a vertex of $ t_i $ by $ \boldsymbol{x} $. Denote by $ t_i^j $ the triangle obtained from $ t_i $ by replacing $ \boldsymbol{v}^{(i)}_j $ with $ \boldsymbol{x} $. Then, $ \boldsymbol{x} \in t_i $ if and only if $ \left| t_i \right| = \left| t_i^1 \right| + \left| t_i^2 \right| + \left| t_i^3 \right| $. Notice that this only requires the computation of vector products, or equivalently determinants, in place of the resolution of a linear system.

This condition alone does not guarantee that the mapping $  \boldsymbol{x} \mapsto t_i $ is well-defined. Indeed, suppose that $ \boldsymbol{x} $ belongs to an edge which is common to two triangles. Then, the admissibility condition holds for both triangles, and the association $ \boldsymbol{x} \mapsto t_i $ does not constitute a function. The  situation gets worse if $ \boldsymbol{x} = \boldsymbol{v}_j^{(i)} $ for some $i$ and $ j \in \{1,2,3\}$; then $ \boldsymbol{x} $ may belong to an even larger number of triangles in $ \mathcal{T}_N$. In these situations, we resolve possible ambiguities by ordering the triangles in $ \mathcal{T}_N$ and by linking $ \boldsymbol{x} $ with the first admissible triangle in the list. 

\begin{remark}
In numerical experiments we could observe that if the triangulation $\T_N$ bears some regularity, for instance, by considering triangles with similar area and avoiding that too many triangles share the same vertex, also different attribution rules $ \boldsymbol{x}^{\mathrm{Pad}} \mapsto t_i $ yield comparable results in terms of the properties of the approximation polynomial.  
\end{remark}



The injectivity of the mapping \eqref{eq:PPtoPT} can be related to the maximal edge length of the triangles in the triangulation $\mathcal{T}_N$. Indeed, let
$$ h_{\max} \doteq \max_{t_i \in \mathcal{T}_N} \max_{\boldsymbol{v}_j^{(i)}, \boldsymbol{v}_k^{(i)} \in t_i} \left\lVert \boldsymbol{v}_j^{(i)} - \boldsymbol{v}_k^{(i)} \right\rVert_2 $$
be the maximum edge length in the triangulation $ \mathcal{T}_N$. If this length is less than the distance between any two Padua points, the map $ \varphi $ is injective. The explicit representation \eqref{PaduaPoints} of the Padua points makes it possible to link the polynomial degree $ m $ with the length $ h_{\max} $.

\begin{proposition}\label{thmimp}
Let $\mathcal{T}_N$ be a triangulation of $\Omega=[-1,1]^2$. If 
\begin{equation}\label{gencond}    m<\frac{\pi}{\arccos\left(1-\sqrt{\frac{h^2_{\max}}{2}}\right)}-1,
    \end{equation}
    each Padua triangle in $\mathcal{T}_M^{\mathrm{Pad}}$ contains exactly one Padua point. Equivalently, the map $\varphi$ in \eqref{eq:PPtoPT} is injective.
\end{proposition}
\begin{proof}
Let $m \in \mathbb{N}$, and consider the set of Padua points $X_M^{\mathrm{Pad}}$ defined in \eqref{PaduaPoints}. To obtain injectivity of the map $\varphi$, we must guarantee that the distance between two Padua points is less than $h_{\max}$. Observing that the Padua points accumulate in the corners of $ \Omega $ (see also \cite{Bos;2006:BLI}), we get the inequalities
\begin{eqnarray*}
\min_{\substack{\boldsymbol{x}, \boldsymbol{y} \in X_M^{\mathrm{Pad}} \\ \boldsymbol{x} \ne \boldsymbol{y} }} \left\lVert \boldsymbol{x} - \boldsymbol{y} \right\rVert_2^2 &\ge& \left( 1 - \cos\left( \frac{\pi}{m+1} \right) \right)^2+\left( 1 - \cos\left( \frac{\pi}{m} \right) \right)^2\\ &>&2\left( 1 - \cos\left( \frac{\pi}{m+1} \right) \right)^2 > h_{\max}^2 .
\end{eqnarray*}
By solving the last inequality for $ m $, we obtain the result.
\end{proof}

As a computational consequence of this result, if the triangulation $ \mathcal{T}_N $ satisfies the hypotheses of Proposition \ref{thmimp}, the following Algorithm \ref{algPad} returns a set of Padua triangles $ \mathcal{T}^{\mathrm{Pad}}_M $ with the correct size $M$ corresponding to the dimension of $\P_m(\mathbb{R}^2)$.

\begin{algorithm}[h!]
	\caption{Extraction of Padua triangles}
	\begin{algorithmic}[1]
		\Require a triangulation $\mathcal{T}_N =\left\{ t_1,\dots,t_N \right\} $ of $ \Omega $, the Padua points $ X_M^{\mathrm{Pad}} $	
		\For{$ \boldsymbol{x}^{\mathrm{Pad}} \in X_M^{\mathrm{Pad}} $}
			\For{ $ i \in \{1, \ldots, N\} $}
                \State check if $\boldsymbol{x}^{\mathrm{Pad}} \in t_i$ by the point-in-triangle algorithm
				\If {$ \boldsymbol{x}^{\mathrm{Pad}} \in t_i $}
				\State add $ t_i $ to $\mathcal{T}^{\mathrm{Pad}}_M$
				\State break inner loop
				\EndIf
			\EndFor
			\EndFor
    \Ensure the set of Padua triangles $\mathcal{T}^{\mathrm{Pad}}_M$
	\end{algorithmic}
	\label{algPad}
\end{algorithm}

\subsection{Regular Friedrichs–Keller triangulations}
Of course, the maximal distance $ h_{\max} $ depends on the chosen triangulation $ \mathcal{T}_N$. For particular regular triangulations, the maximal distance $ h_{\max} $ can be estimated explicitly. As an example, we consider a regular Delaunay triangulation $\mathcal{T}^{\mathrm{reg}}_{N}$ constructed upon the quadratic grid of vertices 
$$ \left( \frac{2 i }{n} - 1, \frac{2 j}{n} - 1 \right) , \quad i,j \in \{0, \ldots, n\}.$$
Specifically, the triangulation $\mathcal{T}^{\mathrm{reg}}_{N}$ is obtained by connecting each interior node of the quadratic grid to six of its eight neighbors, as illustrated in Figure \ref{fig:PadTriangle}. It forms a simple Delaunay triangulation \cite[Chapter 1]{Edelsbrunner} which is also known as Friedrichs–Keller triangulation \cite[p. 64]{Knabner}. This triangulation finds large application in numerical analysis, particularly in finite element methods \cite{Rahla}. It is easy to check that $ \mathcal{T}^{\mathrm{reg}}_{N} $ contains $ N = 2n^2 $ triangles. In Figure~\ref{fig:PadTriangle}, along with the triangulation $ \mathcal{T}^{\mathrm{reg}}_{N}$, also the Padua points $X_M^{\mathrm{Pad}} $ and the triangles $ \mathcal{T}_M^{\mathrm{Pad}} $ extracted by Algorithm \ref{algPad} are shown.  

\begin{figure}[ht]
    \centering
    \includegraphics[width=0.49\textwidth]{TriangPadPnts.eps}
        \includegraphics[width=0.49\textwidth]{TriangPadPntscolor.eps}
    \caption{Padua points (\textcolor{red}{*}) and the corresponding Padua triangles for $m = 6$ in $\mathcal{T}^{\mathrm{reg}}_{800}$ (i.e. with $n = 20$). }
    \label{fig:PadTriangle}
\end{figure}
%
%
%In the hypothesis that $m$ satisfies~\eqref{gencond}, we can define the set of \textit{Padua triangles} as the subset of the triangulation $\mathcal{T}_{N}$ formed by the triangles containing a Padua node, that is
%\begin{equation*}
%\mathcal{T}_{\kappa(m)}^{\mathrm{Pad}}=\left\{t^{\mathrm{Pad}}_{\iota}\, :\,   \iota=1,\dots,\kappa(m)\right\}\subset \mathcal{T}_{N}, 
%\end{equation*}
%such that
%\begin{equation*}
%\boldsymbol{x}^{\mathrm{Pad}}_{\alpha_\iota}\in t_{{\iota}}^{\mathrm{Pad}}, \qquad \iota=1,\dots,\kappa(m).
%\end{equation*}
%

\vspace{2mm}

The particular construction of the Friedrichs-Keller triangulation allows to determine an exact value for $ h_{\max} $ as a function of the step size $2/n$, and hence to estimate the required number of triangles $N$ to get the injectivity of the selection map, making it possible to apply Eq. \eqref{gencond} directly to this particular triangulation.

\begin{corollary}
Let $\Omega=[-1,1]^2$ be partitioned by the regular Friedrichs–Keller triangulation $ \mathcal{T}^{\mathrm{reg}}_{N} $ with $N = 2n^2$. Then,
%\begin{equation}\label{triangulation}
%\mathcal{T}^{\mathrm{reg}}_{N}=\left\{t_i\, : \, i=1,\dots,N\right\}, \qquad N=N:= 2(n-1)^2,
%\end{equation}
%based on the regular Cartesian grid of $n\times n$ equispaced points, see Figure~\ref{fig:PadTriangle}. In this case, we get 
\begin{equation*}
    h_{\max} = \frac{2 \sqrt{2}}{n} = \frac{4}{\sqrt{N}},
\end{equation*}
and the condition~\eqref{gencond} of Proposition~\ref{thmimp} can be simplified to
\begin{equation}\label{conditionOnm}
    m<\frac{\pi}{\arccos\left(\frac{n-2}{n}\right)}-1.
\end{equation}
\end{corollary}

%
%
%
%\begin{remark}
%    We observe that if $m$ satisfies~\eqref{gencond}, the set of Padua triangles is well-defined and satisfies 
%    \begin{equation*}
%        \#\left(\mathcal
%{T}_{m}^{\mathrm{Pad}}\right)=\kappa(m).
%    \end{equation*}
% Moreover, by the mean value theorem, interpolating on these triangles is equivalent to interpolating on points belonging to these triangles. Therefore, as the number of the elements of the triangulation $N$ increases, this tends to interpolation on the Padua nodes.
%\end{remark}
%
%
%In the following, we assume that $m$ satisfies~\eqref{gencond} and we set $r\in \mathbb{N}$ such that $r>m$. We denote by
%\begin{equation}\label{kr}
%\kappa(r):=\binom{r+2}{2},
%\end{equation} 
%and we set
%\begin{equation}\label{basis}
%    \mathcal{B}_r=\left\{e_1(\boldsymbol{x}),\dots,e_{\kappa(r)}(\boldsymbol{x})\right\}, \qquad \boldsymbol{x}=(x,y)\in\Omega,
%\end{equation} 
%a basis of the polynomial space $\mathbb{P}_r\left(\mathbb{R}^2\right)$ of degree less than or equal to $r$ such that
%\begin{equation}\label{propimpbas}
%\operatorname{span}\left\{e_1,\dots,e_{\kappa(m)}\right\}=\mathbb{P}_m\left(\mathbb{R}^2\right)
%\end{equation}
%where $\kappa(m)$ is defined in~\eqref{kappam}.
%Moreover, we assume that the triangulation  $\mathcal{T}_{N}$ is reordered so that 
%\begin{equation*}
%t_{\iota}=t^{\mathrm{Pad}}_{\iota}, \qquad i=1,\dots,\kappa(m).
%\end{equation*}
%Under this assumption, we define the linear functionals 
%\begin{equation}\label{linfuns}
%\nu_{\iota}: f(\boldsymbol{x})\in C\left(\Omega\right) \rightarrow \nu_{\iota}(f)=\int_{t_\iota} f(\boldsymbol{x}) d\boldsymbol{x}, \qquad \iota=1,\dots,N
%\end{equation}
%and we denote by 
%\begin{equation}\label{operatorinterpPad}
%\begin{array}{rcl}
%{P}^{\mathrm{Pad}}_{m}: f(\boldsymbol{x})\in C\left(\Omega\right) &\rightarrow& {P}^{\mathrm{Pad}}_{m}[f](\boldsymbol{x})=\sum\limits_{i=1}^{\kappa(m)}a_ie_i(\boldsymbol{x})\in \mathbb{P}_m\left(\mathbb{R}^2\right), \qquad \boldsymbol{x}\in \Omega, 
%\end{array}
%\end{equation}
%where the vector of coefficients $[a_1,\dots,a_{\kappa(m)}]^T$ is computed such that 
%\begin{equation*}
%\nu_i(f)=\nu_i\left({P}^{\mathrm{Pad}}_{m}\right), \qquad i=1,\dots,\kappa(m).
%\end{equation*}
%The operator ${P}^{\mathrm{Pad}}_{m}$ is also called polynomial histopolation operator on the Padua triangles $\mathcal{T}_M^{\mathrm{Pad}}$.

\subsection{Histopolation on the Padua triangles}

The polynomial interpolation problem in $\P_m\left(\mathbb{R}^2\right)$ based on the Padua nodes $ X_M^{\mathrm{Pad}} $ is uniquely solvable and well-conditioned. In a more formalized way, the Padua points form a \emph{norming set} \cite{NormingSet} for the space $ \P_m\left(\mathbb{R}^2\right)$ with a norming constant $ \mathscr{L}_m \left(X_M^{\mathrm{Pad}}\right)$ that coincides with the (nodal) Lebesgue constant of the polynomial interpolation on the Padua points. Further, unisolvence of norming sets is \emph{stable}, in the sense that a slight perturbation of the nodes $ X_M^{\mathrm{Pad}} $ (or more generally also for sets of averaging functionals \cite{MockSegments}) is still unisolvent. The maximal possible perturbation is estimated in \cite[Proposition 1]{PiazzonVianello} and, for $ \Omega = [-1, 1]^2$, it amounts to a perturbation of up to $ \varepsilon = \frac{\alpha}{2 m^2 \mathscr{L}_m \left(X_M^{\mathrm{Pad}}\right)} $ for any $ \alpha \in [0,1) $.

If the triangulation $ \mathcal{T}_N$ satisfies the hypothesis of Proposition \ref{thmimp}, the cardinality of the Padua triangles $ \mathcal{T}_M^{\mathrm{Pad}} $ coincides with $ M = \dim \left(\P_m\left(\mathbb{R}^2\right)\right)$. The following result gives a necessary condition on $ h_{\max} $ in order to guarantee the unisolvence of the set $ \mathcal{T}_M^{\mathrm{Pad}} $ for histopolation -- or equivalently, the well-posedness of the histopolation operator
\begin{equation} \label{ea:histopolationPadua}
\Pi_{m}^{\mathrm{Pad}} : C (\Omega) \to \P_m\left(\mathbb{R}^2\right)
\end{equation}
that maps a function $f$ to a polynomial $\Pi_{m}^{\mathrm{Pad}}[f]$ of degree $m$ satisfying the histopolation conditions \eqref{eq:histopol} on the Padua triangles $\T_M^{\mathrm{Pad}}$.

\begin{proposition} \label{prop:stability}
	Let $ \mathcal{T}_N $ be a triangulation of $ \Omega $. If for some $0 < \alpha < 1$ we have
	\begin{equation} \label{eq:condhmax}
		h_{\max} < \min \left\{\frac{\alpha}{2 K \left(m \ln m\right)^2}, \sqrt{2} \left( 1 - \cos \left( \frac{\pi}{m+1} \right) \right) \right\}, 
		\end{equation}
where $K$ is linked to the Lebesgue constant of the Padua points, then
	the set $ \mathcal{T}_M^{\mathrm{Pad}} $ is unisolvent in $ \P_m \left(\mathbb{R}^2\right) $ and the histopolation operator $\Pi_{m}^{\mathrm{Pad}}$ well-defined.
\end{proposition}

\begin{proof}
	If the condition \eqref{eq:condhmax} is satisfied, in particular if $ h_{\max} < \sqrt{2} \left( 1 - \cos \left( \frac{\pi}{m+1} \right) \right) $, the attribution map $ X_M^{\mathrm{Pad}} \to \mathcal{T}_M^{\mathrm{Pad}} $ is a bijection by Proposition \ref{thmimp}. Since the mean value theorem \cite[p. 126]{Zorich} states that for a continuous function $f$ we get $ \frac{1}{|t|} \int_{t} f(x,y) \de x \de y = f(\boldsymbol{p}) $ for some point $ \boldsymbol{p} \in t $ (more precisely, in its interior), the histopolation condition on some triangle $ t \in \mathcal{T}_M^{\mathrm{Pad}}$ is equivalent to an interpolation condition on some unknown node $\boldsymbol{p}  \in t$ that depends on $f$. We consider then $ \delta = \left\Vert \boldsymbol{x}^{\mathrm{Pad}} - \boldsymbol{p} \right\Vert_2 $, where $\boldsymbol{x}^{\mathrm{Pad}}$ denotes the Padua point lying in the triangle $t$. Now, since $$ \delta < h_{\max} < \frac{\alpha}{2 K \left(m \ln m\right)^2} \leq \frac{\alpha}{2 m^2 \mathscr{L}_m \left(X_M^{\mathrm{Pad}}\right)},$$ the prerequisites of \cite[Propostion 1]{PiazzonVianello} on the distance between $\boldsymbol{x}^{\mathrm{Pad}}$ and its perturbation $\boldsymbol{p}$ are satisfied; we also used here the fact that the Lebesgue constant for the Padua points satisfies $\mathscr{L}_m \left(X_M^{\mathrm{Pad}}\right) \leq K \left(\ln m\right)^2$ for some constant $K > 0$. Since this bound holds for any triangle $ t \in \mathcal{T}_M^{\mathrm{Pad}}$ and any function $f \in C(\Omega)$, \cite[Propostion 1]{PiazzonVianello} guarantees that the averaging functionals on the Padua triangles $ \mathcal{T}_M^{\mathrm{Pad}} $ provide a norming set for $ \P_m\left(\mathbb{R}^2\right) $ and, hence, the set $ \mathcal{T}_M^{\mathrm{Pad}}$ is unisolvent for histopolation in $ \P_m \left(\mathbb{R}^2\right) $. The estimate $\mathscr{L}_m \left(X_M^{\mathrm{Pad}}\right) \leq K \left(\ln m\right)^2 $ for the Lebesgue constant of the Padua points is, for instance, given in \cite{Bos;2006:BLI}. 
\end{proof}



\subsection{Histopolation-regression based on the Padua triangles} \label{sect:regressionmethod}

It may happen that the data values $\mu_i(f)$ originate from an already collected set of samplings that contains also average values over non-Padua triangles, i.e., averages of $f$ over triangles in the set $ \mathcal{T}_N \setminus \mathcal{T}_M^{\mathrm{Pad}} $. In place of discarding such samplings, one may perform an additional regression on such data to improve the quality of the approximant. For this reason, we formalize the histopolation-regression problem described in \eqref{consleastsquares} mathematically as a constrained least-squares problem~\cite{DellAccio:2022:GOT,dell2024extension,de2024mixed}.  

%Let $f$ be an unknown function and we assume to know only the data~\eqref{eq:mui} and let$ M $ be as defined in~\eqref{M}. We set 

To this end, we consider an integer $ d > m $ such that 
\begin{equation*}\label{DeM}
D-M \doteq \dim\left(\P_d\left(\mathbb{R}^2\right)\right) - \dim \left(\P_m\left(\mathbb{R}^2\right) \right) \leq \# \mathcal{T}_N - \# \mathcal{T}_M.    
\end{equation*}
We assume that $ \{ p_1, \ldots, p_D \} $ forms a basis for $\P_d \left(\mathbb{R}^2\right) $ and is ordered such that $ \{ p_1, \ldots, p_M \} $ spans $ \P_m \left(\mathbb{R}^2\right) $. 
We also assume that the set $\T_N$ is ordered so that its first $M$ elements are the Padua triangles $\T_M^{\mathrm{Pad}}$. Under these assumptions, we define the matrices 
\begin{eqnarray} \label{eq:rectVanderm}
W&\doteq&\left[W_{i,j}\right]=\left[\mu_i\left(p_j\right)\right]\in\mathbb{R}^{N \times D}, \\ C&\doteq&\left[C_{i,j}\right]=\left[\mu_i\left(p_j\right)\right]\in\mathbb{R}^{M \times D},
\end{eqnarray}
and the data vectors 
\begin{equation*}\label{bandd}
\boldsymbol{b}=\left[b_1,\dots,b_N\right]^T \doteq \left[\mu_1(f),\dots,\mu_{N}(f)\right]^T, \quad \boldsymbol{d}=\left[d_1,\dots,d_M\right]^T \doteq \left[\mu_1(f),\dots,\mu_M (f)\right]^T.
\end{equation*}

Using this terminology, we formulate the histopolation-regression problem~\eqref{consleastsquares} mathematically 
as the constrained least-squares problem
\begin{equation} \label{LSProblem} \min_{\boldsymbol{a}\in \mathbb{R}^{D}}\left\lVert W\boldsymbol{a}-\boldsymbol{b}\right\rVert_2, \quad \text{subject to } \quad C \boldsymbol{a} = \boldsymbol{d},
\end{equation} 
where $\left\lVert \cdot\right\rVert_2$ denotes the discrete $2$-norm. As the first $M$ elements of $\T_N$ correspond to the Padua triangles, the matrix $C$ is of full rank $M$ under the assumptions of Proposition \ref{prop:stability}. Further, if we assume that also the matrix $W$ is of full rank, the minimization problem \eqref{LSProblem} has a unique solution. The latter can, for example, be guaranteed if $\T_N$ contains the Padua triangles $\T_D^{\mathrm{Pad}}$ and the respective assumptions of Proposition \ref{prop:stability} are satisfied. 

The solution of the constrained minimization problem~\eqref{LSProblem} can be derived using the method of Lagrange multipliers, as for instance described in \cite{boyd2018introduction}. For this, we require the Lagrangian function
\begin{equation*}
    \mathcal{F}(\boldsymbol{a},\boldsymbol{z})=\left\lVert W\boldsymbol{a}-\boldsymbol{b}\right\rVert_2^2+\sum_{k=1}^Mz_k\left(\boldsymbol{c}_k^T\boldsymbol{a}-d_k\right), \quad \boldsymbol{a}\in\mathbb{R}^D, \, \boldsymbol{z}\in\mathbb{R}^M,
\end{equation*} 
where $\boldsymbol{c}_i^T$ denotes the $i$-th row of the matrix $C$.
If $\hat{\boldsymbol{a}}$ is a solution of~\eqref{LSProblem}, then there exists a vector $\hat{\boldsymbol{z}}=\left[\hat{z}_1,\dots,\hat{z}_{M}\right]^T$ of Lagrange multipliers such that the following conditions hold:
\begin{eqnarray*}
    {\partial_{a_i} \mathcal{F}}\left(\hat{\boldsymbol{a}},\hat{\boldsymbol{z}}\right)&=&0, \quad i=1,\dots,D,\\
    {\partial_{z_j} \mathcal{F}}\left(\hat{\boldsymbol{a}},\hat{\boldsymbol{z}}\right)&=&0,\quad j=1,\dots,M.
\end{eqnarray*}
These $D + M$ conditions can equivalently be expressed in matrix form as
\begin{equation}\label{LagrangeMultmethod}
	\begin{bmatrix}
		2 W^T W & C^T  \\
		C & 0  \\
	\end{bmatrix}
	\begin{bmatrix}
		\hat{\boldsymbol{a}} \\
		\hat{\boldsymbol{z}} \\
	\end{bmatrix}=
	\begin{bmatrix}
		2 W^T\boldsymbol{b} \\
		\boldsymbol{d} \\
	\end{bmatrix}.
\end{equation}
% where
% \begin{equation*}\label{bandd}
% 	\boldsymbol{b} \doteq \left[\mu_1(f),\dots,\mu_{N}(f)\right]^T, \quad \boldsymbol{d} \doteq \left[\mu_1(f),\dots,\mu_M (f)\right]^T,
% \end{equation*}
% and for $ j = 1, \ldots, D $, one has
% \begin{eqnarray}\label{eq:rectVanderm}
%     W_{i,j} &\doteq& \mu_i\left(p_j\right), \quad i=1,\dots,N, \quad j=1,\dots,D,\\
%     C_{i,j} &\doteq& \mu_i\left(p_j\right), \quad i=1,\dots,M, \quad j=1,\dots,D. \nonumber
% \end{eqnarray}

Once this system is solved, the first $ D $ elements of the solution  $\hat{\boldsymbol{a}}=\left[\hat{a}_1,\dots,\hat{a}_D \right]^T$ allow to define the histopolation-regression operator
\begin{align}\label{operatorPhat}
\hat{\Pi}_d^{\mathrm{Pad}} : \quad C (\Omega) & \to \P_d\left(\mathbb{R}^2\right), \\ \nonumber
	f & \mapsto \sum_{i=1}^D \hat{a}_i p_i,
%\begin{array}{rcl}
%\hat{P}^{\mathrm{Pad}}_{r}: f(\boldsymbol{x})\in C\left(\Omega\right) &\rightarrow& \hat{P}^{\mathrm{Pad}}_{r}[f](\boldsymbol{x})=\sum\limits_{i=1}^{\kappa(r)}\hat{a}_ie_i(\boldsymbol{x})\in \mathbb{P}_r\left(\mathbb{R}^2\right), \qquad \boldsymbol{x}\in \Omega, 
%\end{array}
\end{align}
that maps a continuous function to the polynomial approximant $\hat{\Pi}_d^{\mathrm{Pad}}[f]$.

\begin{remark}
   There are several scenarios in which the usage of a histopolation-regression approach \eqref{LSProblem} has advantages compared to a pure histopolation method. For instance, it might happen that the number of triangles $N$ in the triangulation $\mathcal{T}_N$ does not exactly match the dimension $\binom{m+2}{2}$ of the space of polynomials $ \P_m\left(\mathbb{R}^2\right) $. In this case, while the construction of a histopolant requires the selection of an appropriate subset of $\mathcal{T}_N$, the histopolation-regression approach can exploit the entire data on $\mathcal{T}_N$. Furthermore, similar to the nodal case, increasing the number of triangles used for histopolation leads to an increasingly bad-conditioned histopolation problem, yielding larger error propagations in the calculation of the solutions. The histopolation-regression strategy is able to reduce this bad-conditioning if the parameters $m$ and $d$ are appropriately selected.
\end{remark}




%where $\hat{\boldsymbol{a}}=\left[\hat{a}_1,\dots,\hat{a}_{\kappa(r)}\right]^T$ is the solution of the KKT linear equations~\cite{boyd2018introduction}
%\begin{equation}\label{LagrangeMultmethod}
%    \begin{bmatrix}
%2 V^TV & C^T  \\
%C & 0  \\
%\end{bmatrix}
%\begin{bmatrix}
%\hat{\boldsymbol{a}} \\
%\hat{\boldsymbol{z}} \\
%\end{bmatrix}=
%\begin{bmatrix}
%2 V^T\boldsymbol{b} \\
%\boldsymbol{d} \\
%\end{bmatrix}
%\end{equation}
%with 
%\begin{equation}\label{matrixV}
%  V=\begin{bmatrix}
%\nu_1\left(e_1\right) & \nu_1\left(e_2\right) & \cdots & \nu_{1}\left(e_{\kappa(r)}\right)\\
%\nu_2\left(e_1\right) & \nu_2\left(e_2\right) & \cdots & \nu_{2}\left(e_{\kappa(r)}\right)\\
%\vdots  & \vdots  & \ddots & \vdots  \\
%\nu_{N}\left(e_1\right) & \nu_{N}\left(e_2\right) & \cdots & \nu_{N}\left(e_{\kappa(r)}\right)\\
%\end{bmatrix},
%\end{equation}
%\begin{equation}\label{matrixC}
%    C=\begin{bmatrix}
%\nu_1(e_1) & \nu_1(e_2) & \cdots & \nu_{1}(e_{\kappa(r)})\\
%\nu_2(e_1) & \nu_2(e_2) & \cdots & \nu_{2}(e_{\kappa(r)})\\
%\vdots  & \vdots  & \ddots & \vdots  \\
%\nu_{\kappa(m)}(e_{1}) & \nu_{\kappa(m)}(e_2) & \cdots & \nu_{\kappa(m)}(e_{\kappa(r)})\\
%\end{bmatrix},
%\end{equation}
%
%\begin{equation}\label{bandd}
%    \boldsymbol{b}=[\nu_1(f),\dots,\nu_{N}(f)]^T, \qquad  \boldsymbol{d}=[\nu_1(f),\dots,\nu_{\kappa(m)}(f)]^T
%\end{equation} 
%and $\hat{\boldsymbol{z}}$ is the Lagrange multipliers vector.  

\begin{remark}
    By construction, the histopolation-regression operator \eqref{operatorPhat} is a linear projector, that is
    $$ \hat{\Pi}_d^{\mathrm{Pad}} \left[\lambda_1 f + \lambda_2 g\right] = \lambda_1 \hat{\Pi}_d^{\mathrm{Pad}}[f] + \lambda_2 \, \hat{\Pi}_d^{\mathrm{Pad}}[g], \quad \forall \ \lambda_1,  \lambda_2 \in \R, \ \ f,g \in C(\Omega), $$
    and 
    \begin{equation}\label{cind2}
     \hat{\Pi}^{\mathrm{Pad}}_{d}\left[\hat{\Pi}^{\mathrm{Pad}}_{d} [f]\right]=\hat{\Pi}^{\mathrm{Pad}}_{d} [f], \qquad f \in C(\Omega).
    \end{equation}
%    \begin{equation*}\label{cind1}
%        \hat{P}^{\mathrm{Pad}}_{r}[f+\lambda g]=\hat{P}^{\mathrm{Pad}}_{r}[f]+\lambda\hat{P}^{\mathrm{Pad}}_{r}[g], 
%    \end{equation*}
%    for any $f,g\in C(\Omega)$ and $\lambda\in\mathbb{R}$. Moreover, it satisfies
%    \begin{equation}\label{cind2}
%        \hat{P}^{\mathrm{Pad}}_{r}[p]=p, \qquad p\in \mathbb{P}_r(\mathbb{R}^2).
%    \end{equation}
As a consequence, for each $ f \in C (\Omega) $, one has a Lebesgue-type inequality of the form
\begin{equation} \label{eq:lebesguelikeinequality}
	 \left\Vert f -  \hat{\Pi}_d^{\mathrm{Pad}}[f] \right\Vert_\infty \leq \left( 1 + \left\lVert \hat{\Pi}_d^{\mathrm{Pad}} \right\rVert_{\mathrm{op}} \right) \min_{p \in \P_d\left(\mathbb{R}^2\right)} \left\Vert f - p \right\Vert_\infty .
\end{equation}
\end{remark}
%QUI E PRIMA SERVONO UN SACCO DI REFERENZE.

\subsection{Bounding the Lebesgue constant for the histopolation-regression method}

A common strategy to quantify the numerical conditioning of an approximation or interpolation operator $ \Pi $, is to consider its operator norm $ \left\Vert \Pi \right\Vert_{\mathrm{op}} $. Such an operator norm is induced by a vector space norm for the corresponding function spaces, and in many cases this norm corresponds to the uniform norm. In several contexts, the quantity $ \left\Vert \Pi \right\Vert_{\mathrm{op}} $ is termed \emph{Lebesgue constant}, and it is readily seen in the Lebesgue inequality \eqref{eq:lebesguelikeinequality} that it quantifies the error between a function and its polynomial approximant in terms of the best polynomial approximation. More detailed explanations on the role of the Lebesgue constant can be found in \cite[Chapter 2]{Davis:1975:IAA}.

To get an accessible bound of the operator norm $ \| \hat{\Pi}_d^{\mathrm{Pad}} \|_{\mathrm{op}} $ in case of the histopolation-regression operator $\hat{\Pi}_d^{\mathrm{Pad}}$, we follow the direct elimination method in \cite[Chapter 5]{bjorck2024numerical}, and have a closer look at the computation of the coefficients $\hat{\boldsymbol{a}}$ in the expansion \eqref{operatorPhat}. The idea consists of splitting $\hat{\boldsymbol{a}} =\left[\hat{\boldsymbol{a}}_1, \hat{\boldsymbol{a}}_2\right]^T $ into two parts, with $ \hat{\boldsymbol{a}}_1 \in \R^M $ and $ \hat{\boldsymbol{a}}_2 \in \R^{D-M} $. Then, we may factorize $ C = QR = Q \left[R_1 , R_2 \right]$, being $ R_1 \in \R^{M \times M} $ a nonsingular, upper triangular matrix. %In what follows, it is convenient to think of $ R $ as the juxtaposition of two matrices $ R = [R_1, R_2] $ 
%\begin{equation}\label{vectdec}
%\hat{\boldsymbol{a}} 
%=
%\begin{bmatrix}
%\hat{\boldsymbol{a}}_1 \\
%\hat{\boldsymbol{a}}_2
%\end{bmatrix}, \qquad \hat{\boldsymbol{a}}_1 \in \mathbb{R}^{\kappa(m)}, \qquad \hat{\boldsymbol{a}}_2 \in \mathbb{R}^{\kappa(r)-\kappa(m)}.
%\end{equation}
%More in details, we consider the QR factorization of the matrix $C$, that is
%\begin{equation*}
%    C = Q_C R_C,
%\end{equation*}
%where $Q_C \in \mathbb{R}^{\kappa(m)\times \kappa(m)}$ is orthogonal, $R_C=[R_{11},R_{12}]\in \mathbb{R}^{\kappa(m)\times \kappa(r)}$ and 
%\begin{equation*}
%  R_{11} \in \mathbb{R}^{\kappa(m) \times \kappa(m)}
%\end{equation*}
%is a nonsingular upper triangular matrix. 
%In this setting, it can be shown that
Plugging this into \eqref{LagrangeMultmethod}, one immediately retrieves
\begin{equation}\label{a1eq}
\hat{\boldsymbol{a}}_1 = R_{1}^{-1} \left(Q^T\boldsymbol{d} - R_{2} \hat{\boldsymbol{a}}_2\right).    
\end{equation}
Likewise, splitting $ W = \left[W_1 , W_2\right] $ with $ W_1 \in \R^{N \times M} $, one obtains
\begin{equation}\label{a2eq}
	\hat{\boldsymbol{a}}_2=\left(A^{T} A\right)^{-1}A^{T}\boldsymbol{b}_1,
\end{equation}
where we have set $ A \doteq W_2 - W_1 R_{1}^{-1} R_{2} $ and $ \boldsymbol{b}_1 \doteq \boldsymbol{b} - W_1 R_{1}^{-1} Q^T\boldsymbol{d} $. We follow now the strategy given in \cite{Submitted:2025:OTE} to bound $ \| \hat{\Pi}_d^{\mathrm{Pad}} \|_{\mathrm{op}} $ in terms of these factors, isolating and identifying the contributions from $ \hat{\boldsymbol{a}}_1 $ and $ \hat{\boldsymbol{a}}_2 $.
%To compute $\hat{\boldsymbol{a}}_2$, we first decompose the matrix $V$ as
%\begin{equation*}
%  V = [V_1, V_2], \qquad V_1\in \mathbb{R}^{N \times \kappa(m)}.  
%\end{equation*}
%Then, the vector $\hat{\boldsymbol{a}}_2$ can be expressed as
%\begin{equation}\label{a2eq}
%    \hat{\boldsymbol{a}}_2=(M_1^{T}M_1)^{-1}M_1^{T}\boldsymbol{b}_1,
%\end{equation}
%where 
%\begin{equation}\label{M1}
%    M_1=V_2-V_1R_{11}^{-1}R_{12}
%\end{equation}
%and
%\begin{equation}\label{B1}
%\boldsymbol{b}_1 = \boldsymbol{b} - V_1 R_{11}^{-1} Q_{C}^T\boldsymbol{d}. 
%\end{equation}
%For further details, see~\cite[Chapter 5]{bjorck2024numerical}.
%%\end{remark}
%
%Following the same strategy used in~\cite{Submitted:2025:OTE}, we aim to estimate the norm of the Padua histopolation-regression operator $\hat{P}^{\mathrm{Pad}}_{r}$ using the direct elimination method. By equipping $C\left(\Omega\right)$ and $\mathbb{P}_{r}(\mathbb{R}^2)$ with the $L_{\infty}$-norm 
%\begin{equation*}
%    \left\lVert f\right\rVert_{\infty}=\max_{\boldsymbol{x}\in \Omega}\left\lvert f(\boldsymbol{x})\right\rvert, 
%\end{equation*}
% we aim to determine a bound for the norm
%\begin{equation}
%\label{eq:NormOfOperator}
%\left\lVert \hat{P}^{\mathrm{Pad}}_{r} \right\rVert_{\infty}:= \sup\limits_{\substack{f\in C\left(\Omega\right)\\ \left\lVert f \right\rVert_{\infty}\le1}} \left\lVert \hat{P}^{\mathrm{Pad}}_{r}[f] \right\rVert_{\infty}.
%\end{equation} 


\begin{proposition} \label{Thm:EstimateOfNorm}
Let $ \left\{ p_1, \ldots, p_D \right\} $ be a basis for $ \P_d\left(\mathbb{R}^2\right) $ with $\underset{i=1,\dots,D}{\max} \left\lVert p_i \right\rVert_{\infty} \leq 1$. Then,
\begin{equation}
\label{eq:BoundThm}
    \left\lVert \hat{\Pi}_d^{\mathrm{Pad}} \right\rVert_{\mathrm{op}} \leq \zeta_d + \eta_d,
\end{equation}
where
\begin{equation*}
\zeta_d \doteq \left\lVert R_{1}^{-1} \right\rVert_{1} \left(M \left\lVert Q^{T} \right\rVert_{1} + \left\lVert R_{2}\right\rVert_{1}\eta_d\right), \quad     \eta_d \doteq \left\lVert(A^T A)^{-1} A^T \right\rVert_1\left(D + M \left\lVert  W_1 R_{1}^{-1}Q^T\right\rVert_1  \right).
\end{equation*}
\end{proposition}


\begin{proof}
%Let $f\in C(\Omega)$ such that
%\begin{equation*}
%        \left\lVert f\right\rVert_{\infty}=\max_{\boldsymbol{x}\in \Omega}\left\lvert f(\boldsymbol{x})\right\rvert.
%\end{equation*}
%	We observe that
For any continuous function $f\in C(\Omega)$, we have
	\begin{equation} \label{eq:comoda}
%		\left\lVert \hat{\Pi}_d^{\mathrm{Pad}} \right\rVert_{\mathrm{op}} = \sup_{\Vert f \Vert_\infty = 1} 
\left\lVert \hat{\Pi}_d^{\mathrm{Pad}} [f] \right\rVert_\infty = \left\lVert \sum_{i=1}^D \hat{{a}}_i (f) p_i \right\rVert_\infty \leq \sum_{i=1}^D \left\lvert \hat{{a}}_i (f) \right\rvert = \left\Vert \hat{\boldsymbol{a}} (f) \right\Vert_1 = \left\Vert \hat{\boldsymbol{a}}_1 (f) \right\Vert_1 + \left\Vert \hat{\boldsymbol{a}}_2 (f) \right\Vert_1.
		\end{equation}
%Let $f\in C\left(\Omega\right)$ be such that   
%\begin{equation} \label{hyp}
%\left\lVert f \right\rVert_{\infty}\le 1.
%\end{equation}
% Let $\boldsymbol{x}\in\Omega$ and by using the triangular inequality, we obtain
%\begin{equation}
%\label{eq:ContinuityCMCLSO}
%    \left\lvert \hat{P}^{\mathrm{Pad}}_{r}[f](\boldsymbol{x}) \right\rvert=\left\lvert \sum_{i=1}^{\kappa(r)} \hat{a}_i(f)e_i(\boldsymbol{x}) \right\rvert\le \sum_{i=1}^{\kappa(r)} \left\lvert\hat{a}_i(f)\right\rvert \left\lvert e_i(\boldsymbol{x})\right\rvert.
%\end{equation}
%Then, by passing to the maximum to both sides, we have 
%\begin{equation} 
%\label{eq:BoundP(f)}
%   \left \lVert \hat{P}^{\mathrm{Pad}}_{r}[f] \right\rVert_{\infty} \le K\sum_{i=1}^{\kappa(r)} \left\lvert\hat{a}_i(f)\right\rvert=K\left\lVert\hat{\boldsymbol{a}}(f)\right\rVert_1,
%\end{equation}
%where $K$ is defined in~\eqref{eq:ConstantC}. Now, it remains to bound the discrete $L_1$-norm of   $\hat{\boldsymbol{a}}=\hat{\boldsymbol{a}}(f)$. We compute this vector by using the direct elimination method. Then, by~\eqref{vectdec}, we obtain
%\begin{eqnarray*}
%\left\lVert
%\hat{\boldsymbol{a}} 
%\right\rVert_1
%&=& 
%\left\lVert
%\begin{bmatrix}
%\hat{\boldsymbol{a}}_1 \\
%\hat{\boldsymbol{a}}_2
%\end{bmatrix}
%\right\rVert_1
%= \lVert \hat{\boldsymbol{a}}_1 \rVert_1 + \lVert \hat{\boldsymbol{a}}_2 \rVert_1.
%\end{eqnarray*}
Using the relations~\eqref{a1eq} and~\eqref{a2eq}, and applying the triangular inequality, we have
\begin{equation*}
    \left\lVert \hat{\boldsymbol{a}}_1 (f) \right\rVert_1\le \left\lVert R_{1}^{-1}\right\rVert_{1} \left(\left\lVert Q^T \right\rVert_1 \left\lVert \boldsymbol{d} (f) \right\rVert_1 + \left\lVert R_{2}\right\rVert_1 \left\lVert \hat{\boldsymbol{a}}_2 (f)\right\rVert_1\right) 
\end{equation*}
and
\begin{equation*}\label{a2}
   \left\lVert \hat{\boldsymbol{a}}_2 (f)\right\rVert_1\le\left\lVert\left(A^{T}A\right)^{-1}A^{T}\right\rVert_1 \left\lVert \boldsymbol{b}_1 (f) \right\rVert_1.
\end{equation*}
It remains to bound $\left\lVert \boldsymbol{b} (f) \right\rVert_1$ and $\left\lVert \boldsymbol{d} (f) \right\rVert_1$. By the mean value theorem, %Since $ \mathcal{T}_N$ is a partition of $ \Omega $, 
one has
\begin{equation*}\label{aus1}
    \left\lVert \boldsymbol{b} (f) \right\rVert_1=\sum_{i=1}^D \left\lvert \mu_i(f)\right\rvert\le\sum_{i=1} ^D \frac{1}{|t_i|}\int_{t_i}\left\lvert f(x,y)\right\rvert \de x \de y \le \left\Vert f \right\Vert_\infty \sum_{i=1} ^D 1 = D \left\Vert f \right\Vert_\infty
\end{equation*}
and 
\begin{equation*}
\left\lVert \boldsymbol{d} (f) \right\rVert_1=\sum_{i=1}^{M} \left\lvert \mu_i(f)\right\rvert\le \sum_{i=1}^M \frac{1}{|t_i|} \int_{t_i}\left\lvert f(x,y)\right\rvert \de x \de y \le \left\Vert f \right\Vert_\infty \sum_{i=1} ^{M} 1 = M \left\Vert f \right\Vert_\infty.
\end{equation*}
%Finally, by substituting~\eqref{B1} into~\eqref{a2}, applying the triangle inequality, and using~\eqref{aus1}, we obtain from~\eqref{eq:BoundP(f)}
%that
Since
\begin{equation*}
    \left\lVert \hat{\Pi}_d^{\mathrm{Pad}} \right\rVert_{\mathrm{op}} = \sup_{\left\Vert f \right\Vert_\infty \leq 1} \left\lVert \hat{\Pi}_d^{\mathrm{Pad}} [f] \right\rVert_\infty, 
\end{equation*}
plugging the above expressions into the bounds for $ \left\lVert \hat{\boldsymbol{a}}_1 (f) \right\rVert_1 $ and $ \left\lVert \hat{\boldsymbol{a}}_2 (f) \right\rVert_1 $, and taking the supremum over all $f \in C(\Omega)$ with $ \left\Vert f \right\Vert_\infty \leq 1 $, we get $ \left\lVert \hat{\boldsymbol{a}}_1 \right\rVert_1 \leq \zeta_d $ and $ \left\lVert \hat{\boldsymbol{a}}_2\right\rVert_1 \leq \eta_d $. Finally, substituting this into \eqref{eq:comoda}, we get the desired bound in \eqref{eq:BoundThm}.
\end{proof}
%\begin{remark}
%	In Proposition \ref{Thm:EstimateOfNorm} the choice of the Chebyshev basis is made only for convenience, as the norm of its elements on $ \Omega $ is upper bounded by $ 1 $, i.e.
%     \begin{equation*} 
%   \max_{i=1,\dots,D} \left\lVert p_i \right\rVert_{\infty} = 1,
%   \end{equation*}
%    and thus any other basis of norm $1$ functions works. 
%Since the norm of an operator does not depend on its representation, we may conclude that the right hand side of \eqref{eq:BoundThm} provides a bound for the norm of the histopolation-regression operator.% However, it is important to emphasize that the bound of the operator established in \eqref{eq:BoundThm} does not depend on the specific choice of basis. {\color{red} W: da chiarire: penso che dipenda dalla base utilizzata, anche fortemente. Con la base monomiale la stima ci esplode}
%\end{remark}


\begin{figure}[ht]
	\centering
	\includegraphics[width = 0.6\textwidth]{G1G2.eps}
	\caption{Trend of the bound $ \zeta_d + \eta_d $ in \eqref{eq:BoundThm} for increasing numbers $N$ of triangles in the Friedrichs-Keller triangulation $\mathcal{T}_{N}^{\mathrm{reg}}$. This trend is computed using the largest integer $m$ satisfying~\eqref{conditionOnm} and the degree $d = m + \sqrt{m}$.}
	\label{G1G2}
\end{figure}


%
%\begin{remark}
%\label{remark_Norm1ChebPolBasis}
%We notice that, if we express the operator $\hat{P}^{\mathrm{Pad}}_r$ with respect to the basis
%\begin{equation}\label{ChebBasis}
%    \tilde{\mathcal{B}}^C_r=\left\{u^C_i(x)u_j^C(y)\, :\, 0<i+j<r\right\},
%\end{equation}
%where $\{u_i^C\}_{i=0}^r$ is the Chebyshev polynomial basis of the first kind, we obtain
%\begin{equation*}
%    K=1.
%\end{equation*}
%Consequently, the bound~\eqref{eq:BoundThm} becomes 
%\begin{equation}
%\label{eq:newBound}
%    \left\lVert \hat{P}^{\mathrm{Pad}}_{r} \right\rVert_\infty \leq G_1+G_2.
%\end{equation}
%\end{remark}
%



\begin{remark}
Note that the quantities $\zeta_d$ and $\eta_d$ and, thus, the quality of the estimate \eqref{eq:BoundThm} depend strongly on the selected basis. Both, the Chebyshev basis and the monomial basis, satisfy the hypothesis of Proposition \ref{Thm:EstimateOfNorm}. However, while for the monomial basis the matrices $ C $ and $ W $ are highly ill-conditioned, it is known that the Chebyshev basis leads to rather well-conditioned Vandermonde-type matrices. For this, the Chebyshev basis seems to be a good candidate to calculate the bounds $\zeta_d$ and $\eta_d$.
In Figure~\ref{G1G2}, we report the right-hand bounds \eqref{eq:BoundThm} (blue points) for the norm of the histopolation-regression operator for different Friedrichs–Keller triangulations $\mathcal{T}_{N}^{\mathrm{reg}}$, $N = 2n^2$, $n\in\mathbb{N}$, in the square $\Omega=[-1,1]^2 $. The values are computed using the largest integer $m$ satisfying~\eqref{conditionOnm} and the degree $d = m + \sqrt{m}$. 
From this graphic, we observe that 
     \begin{equation*}
     \zeta_d + \eta_d \approx e^{3.64 \log n}=n^{3.64}. 
     \end{equation*}
\end{remark}
%
%As a consequence of Theorem~\ref{Thm:EstimateOfNorm}, we can establish an error bound for the approximation operator $\hat{P}_{r}^{\mathrm{Pad}}$. The following theorem holds
%
%\begin{theorem}
%\label{Thm:Errorb}
%Let $\Omega=[-1,1]^2$ be the domain partitioned by the triangulation $\mathcal{T}_{N}$ and let $f\in C(\Omega)$. Then, by expressing the approximation operator $\hat{P}_{r}^{\mathrm{Pad}}$ in the basis $\tilde{\mathcal{B}}_r^C$, we have
%\begin{equation}   \label{bounderroruniformapprox}
%    \left\lVert f-\hat{P}_{r}^{\mathrm{Pad}}[f]\right\rVert_{\infty}\le \left(1+G_1+G_2\right)\mathcal{E}_{r}(f),
%\end{equation}
%where $\mathcal{E}_{r}(f)$ is the error of best uniform approximation by polynomials of $\mathbb{P}_r\left(\mathbb{R}^2\right)$~\cite{weinstein1972uniform}.
%\end{theorem}
%\begin{proof}
%Let $p^{\star}_r[f]\in \mathbb{P}_r\left(\mathbb{R}^2\right)$ be the polynomial of best uniform approximation to $f$ on the domain $\Omega$. Then, by~\eqref{cind1} and~\eqref{cind2}, we get
%\begin{eqnarray*}
%    \left\lVert f-\hat{P}^{\mathrm{Pad}}_{r}[f]\right\rVert_{\infty}&=&    \left\lVert f-p_r^{\star}[f]+p_r^{\star}[f]-\hat{P}^{\mathrm{Pad}}_{r}[f]\right\rVert_{\infty}\\
%    &=&\left\lVert f-p_r^{\star}[f]-\hat{P}^{\mathrm{Pad}}_{r}[f-p_r^{\star}[f]]\right\rVert_{\infty}\\
%    &\le& \left(1+\left\lVert \hat{P}^{\mathrm{Pad}}_{r}\right\rVert_\infty\right)\left\lVert f-p_r^{\star}[f]\right\rVert_{\infty}.
%\end{eqnarray*}
%The thesis then follows from  Theorem~\ref{Thm:EstimateOfNorm}.
%\end{proof}

\subsection{Bounding the Lebesgue constant in the histopolation setting}
In the pure histopolation setting, we can derive more concrete bounds for the Lebesgue constant, as in this case also the norms $\|\Pi_m^{\mathrm{Pad}}\|_{\mathrm{op}}$ can be characterized more explicitly. These characterizations are well-known in the classical interpolation setting where the Lebesgue constant can be stated in terms of the Lagrange basis functions. Such characterizations have been found also in more abstract frameworks \cite{ARRLeb}. For a collection $ \mathcal{T}_M = \{ t_1, \ldots, t_M \} $ of triangles that form a unisolvent set for $ \P_m \left(\mathbb{R}^2\right)$, we obtain this characterization of the Lebesgue constant in terms of the quantity
\begin{equation} \label{eq:LebConst}
	\mathscr{L}_m \left(\mathcal{T}_M\right) \doteq \sup_{t \subset \Omega} \frac{1}{|t|} \sum_{i=1}^M \left| \int_t \ell_{t_i}(x,y) \de x \de y \right|  = \sup_{\xi \in \Omega} \sum_{i=1}^M \left| \ell_{t_i}(\xi)\right|, 
\end{equation} 
where $  t $ ranges over all triangles supported in $ \Omega $ and $ \{ \ell_{t_1}, \ldots, \ell_{t_M} \} $ is the Lagrange basis for the space $\mathbb{P}_m(\mathbb{R}^2)$ satisfying the histopolation conditions $ \frac{1}{|t_i|} \int_{t_i} \ell_{t_j}(x,y) \de x \de y = \delta_{i,j} $. While the first definition in \eqref{eq:LebConst} corresponds to the definition originally given in \cite{ARRLeb} for differential forms, the second is a consequence of the mean value theorem and resembles the definition of the Lebesgue constant in the nodal setting. The following identity for the Lebesgue constant can be derived from \cite{BruniElefante}.% is given in terms of simplices, the request on the shape of test supports may be relaxed \cite{BruniElefante}; nevertheless, it appears convenient to keep a geometrical consistency with the elements of $ \mathcal{T}_M $, see the discussion in \cite[Section 6]{ABR22}.


\begin{proposition} \label{prop:intnorm}
Let $\mathcal{T}_M = \{ t_1, \ldots, t_M \} $ be a unisolvent set of triangles for $ \P_m \left(\mathbb{R}^2\right)$ and suppose that the triangles in $\mathcal{T}_M$ are disjoint in the sense that the measure $\left| t_i \cap t_j\right| = 0 $ for each $ t_i,t_j \in \mathcal{T}_M $, $t_i\ne t_j$. Then, the operator norm of the respective histopolation operator $\Pi_m$ satisfies
$$ \left\Vert \Pi_m \right\Vert_{\mathrm{op}} = \mathscr{L}_m \left(\mathcal{T}_M\right). $$
    \end{proposition}

In order to calculate the quantity $\mathscr{L}_m \left(\mathcal{T}_M\right)$, and hence also the operator norm $\left\Vert \Pi_m \right\Vert_{\mathrm{op}}$, we require a scheme to compute the Lagrange basis of the space $ \P_m \left(\mathbb{R}^2\right)$ for a unisolvent set $\mathcal{T}_M$ of triangles.
Despite general explicit expressions for the Lagrange basis are not available (see \cite{bruni2024polynomial} for an account), we can calculate the Lagrange basis numerically in terms of a chosen basis $ \{ p_1, \ldots, p_M \} $ of $ \P_m\left(\mathbb{R}^2\right) $. Once this basis is selected, it is sufficient to invert the \emph{generalized Vandermonde matrix}%  
\begin{equation} \label{eq:VandermondeMatrix}
	[V_{i,j}] \doteq [\mu_i(p_j)] = \left[ \frac{1}{|t_i|} \int_{t_i} p_j(x,y) \de x \de y \right] \in \mathbb{R}^{M \times M},
\end{equation}
to recover the expansion coefficients for the Lagrange basis functions. More precisely, we have
$$ \int_{t_i} \sum_{k=1}^M V_{k,j}^{-1} p_k (x,y) \de x \de y = \delta_{i,j} ,$$
meaning that $ \ell_j (x,y) = \sum_{k=1}^M V_{k,j}^{-1} p_k (x,y) $ provides the desired Lagrange basis functions. The Vandermonde matrix \eqref{eq:VandermondeMatrix} is invertible if and only if the functionals $\mu_i(f)$ are linearly independent, see \cite[Lemma 3.2.2]{AtkinsonHan}. In case of the Padua triangles $\T_M^{\mathrm{Pad}}$, the invertibility of the matrix \eqref{eq:VandermondeMatrix} is guaranteed under the assumptions of Proposition \ref{prop:stability}.  


To calculate the Lebesgue constant $\mathscr{L}_m \left(\mathcal{T}_M\right)$ we can now proceed as follows. We consider the rightmost term in the definition \eqref{eq:LebConst} of the constant $ \mathscr{L}_m \left(\mathcal{T}_M\right) $ %and a collection of 
%supports $ \mathcal{T}_N =\left\{s_1, \ldots, s_N \right\}$ (with $ N \gg M $), one finds that $ \mathscr{L}_m \left(\mathcal{T}_M\right) $ may be written as the product of four matrices as
% $$ \mathscr{L}_m \left(\mathcal{T}_M\right) \approx  \left\Vert \mathrm{diag} \left(\left|s_1\right |^{-1}, \ldots, \left|s_N\right |^{-1}\right) W V^{-1} \mathrm{diag} \left( \left|t_1\right|, \ldots, \left|t_M \right|\right) \right\Vert_\infty,$$
% where 
% $$ W_{i,j} = \int_{s_i} p_j (x,y) \de x \de y .$$%= | s_i | p_j(\xi_i) .$$
% By applying the mean value theorem to elements of $ W $, one finds that
% $$ \left|s_i\right|^{-1} W_{i,j} = \left|s_i\right|^{-1} \int_{s_i} p_j (x,y) \de x \de y = p_j\left(\xi_i\right) $$
%for some $ \xi_i \in s_i $. Thus, if $ \mathcal{T}_N $ is large enough, one may consider a comparably large collection of 
%nodes $ \mathcal{X} = \{ \xi_1, \ldots, \xi_N\} $, % and replace $ \mathrm{diag} \left(\left| s_1 \right|^{-1}, \ldots, \left| s_N \right|^{-1}\right) W $ by 
and define the vector
$ \boldsymbol{p} (\xi) \doteq \left[ p_1 \left(\xi\right), \ldots, p_M (\xi) \right] $.
Then, using the inverse of the Vandermonde matrix $V$ to describe the Lagrange basis functions, we can rewrite $\mathscr{L}_m \left(\mathcal{T}_M\right)$ more compactly as
	\begin{equation} \label{eq:shortLeb}
		\mathscr{L}_m \left(\mathcal{T}_M\right) = \sup_{\xi \in \Omega} \left\Vert \boldsymbol{p} (\xi) \, V^{-1}  \right\Vert_1 .
	\end{equation}
%	where %$ V $ is the Vandermonde matrix associated with any basis $ \{ p_1, \ldots, p_m \} $ of $ \mathbb{P}_m\left(\mathbb{R}^2\right)  $ in 
	%and 
	
%    The second one, which is Eq. \eqref{eq:shortLeb}, offers a way to efficiently estimate such a quantity.
% Notice that the estimation of \eqref{eq:LebConst} via \eqref{eq:shortLeb} now requires only quadrature for the Vandermonde matrix \eqref{eq:VandermondeMatrix} and we thus exploit it to estimate the Lebesgue constant \eqref{eq:LebConst} associated with Padua triangles. 
To estimate the Lebesgue constant $\mathscr{L}_m \left(\mathcal{T}_M\right)$ numerically, we discretize $ \Omega $ by a suitable large discrete set of nodes $ \mathcal{X} = \{ \xi_1, \ldots, \xi_N\} $ and then replace $ \Omega $ by $ \mathcal{X} $ in \eqref{eq:shortLeb}. The whole computation can be implemented efficiently by evaluating the infinity norm of a product of two matrices. 
%The above result, which is based on the mean value theorem applied to elements of the matrix $ W $ introduced in Eq. \eqref{eq:rectVanderm}, and simplifies the formula \eqref{eq:LebConst}, avoiding quadrature. In fact, one may expand \eqref{eq:LebConst} 

%We thus exploit \eqref{eq:shortLeb} to estimate the Lebesgue constant \eqref{eq:LebConst} associated with Padua triangles. 
In Figure \ref{fig:LebPD}, the numerically calculated Lebesgue constant associated with the Padua triangles $\mathcal{T}_M^{\mathrm{Pad}}$ is shown. It is visible that the Lebesgue constant $\mathscr{L}_m \left(\mathcal{T}_M^{\mathrm{Pad}} \right)$ increases only slowly for a growing total degree $m$. The following proposition states that this growth is at most logarithmic raised to the power of two if the size of the triangles in $\mathcal{T}_M^{\mathrm{Pad}}$ is sufficiently small.  

\begin{proposition} \label{prop:stability2}  
	If the maximal triangle length $h_{\max}$ is small enough according to the bound in Proposition \ref{prop:stability}, the Lebesgue constant of the Padua triangles can be bounded by
	$$ \mathscr{L}_m \left( \mathcal{T}_M^{\mathrm{Pad}} \right) \leq \frac{K}{1-\alpha} \left(\ln m\right)^2,$$
	where $0 < \alpha < 1$ and $K>0$ is linked to the Lebesgue constant of the Padua nodes.  
\end{proposition}

\begin{proof}
    As in the proof of Proposition \ref{prop:stability}, the mean value theorem implies that the histopolation operator $\Pi_{m}^{\mathrm{Pad}} [f] = \Pi_m(\mathcal{T}_M^{\mathrm{Pad}}) [f]$ providing the polynomial histopolant of the function $f \in C(\Omega)$ on the Padua triangles $\mathcal{T}_M^{\mathrm{Pad}}$ can be read as a nodal interpolation operator $\Pi_m(\Xi_M(f))$$ [f]$ of the function $f$ on points $ \Xi_M(f) = \{ \xi_1(f), \ldots, \xi_M(f) \} $ that depend on the function $f$ and are contained in $\mathcal{T}_M^{\mathrm{Pad}}$, i.e., $ \xi_i(f) \in t_i $ for some Padua triangle $t_i \in \mathcal{T}_M^{\mathrm{Pad}}$. Since the distance between a point $\xi_i(f)$ and the corresponding Padua point in the triangle $t_i$ is smaller than the maximal length $ \max_i \ \mathrm{diam} \left(t_i\right) \leq h_{\max} $, the assumption on $h_{\max}$ ensures that \cite[Corollary 1]{PiazzonVianello} can be applied and that the norm of the interpolation operators associated with $ X_M^{\mathrm{Pad}} $ and $\Xi_M$ satisfy
	\begin{equation} \label{eq:helpme}
	\frac{\left\Vert \Pi_m(\Xi_M(f)) \right\Vert_{\mathrm{op}}}{\left\Vert \Pi_m(X_M^{\mathrm{Pad}}) \right\Vert_{\mathrm{op}}} \leq \frac{1}{1-\alpha}.  
	\end{equation} 
	The equality $ \left\Vert \Pi_m(X_M^{\mathrm{Pad}}) \right\Vert_{\mathrm{op}} = \mathscr{L}_m \left(X_M^{\mathrm{Pad}}\right) $ is classical and follows from the characterization of the nodal Lebesgue constant as the norm of the corresponding nodal interpolator. As the Lebesgue constant $ \mathscr{L}_m \left(X_M^{\mathrm{Pad}}\right) $ for the Padua points can be bounded as $\mathscr{L}_m \left(X_M^{\mathrm{Pad}}\right) \leq K \left(\ln m\right)^2 $, and \eqref{eq:helpme} holds true for any function $f \in C(\Omega)$, we get
    \[ \left\Vert \Pi_m(\mathcal{T}_M^{\mathrm{Pad}}) \right\Vert_{\mathrm{op}} \leq \sup_{ \| f \|_{\infty} \leq 1} \left\Vert \Pi_m(\Xi_M(f)) \right\Vert_{\mathrm{op}} \leq \frac{K}{1-\alpha} \left(\ln m\right)^2. \]
    The identity $ \left\Vert \Pi_m(\mathcal{T}_M^{\mathrm{Pad}}) \right\Vert_{\mathrm{op}} = \mathscr{L}_m \left(\mathcal{T}_M^{\mathrm{Pad}}\right) $ is guaranteed by Proposition \ref{prop:intnorm}, as Padua triangles do not intersect in their interior.
\end{proof}

The trend for the Lebesgue constant $\mathscr{L}_m \left( \mathcal{T}_M^{\mathrm{Pad}} \right) $ is depicted in Figure \ref{fig:LebPD}.
The left subfigure indicates that the Lebesgue constant is primarily affected by the degree $m$ of the polynomial space and only in a minor way on the number $N$ of triangles. Further, for increasing $N$ the Lebesgue constant of the Padua triangles tends to the Lebesgue constant of the Padua points. This gets also visible on the right hand side of Figure \ref{fig:LebPD}, where the trend of the Lebesgue constant of the Padua triangles shows a consistent behavior with that of the Padua points, cf. \cite{Bos;2006:BLI}.

\begin{figure}[ht]
	\centering
    \includegraphics[width=0.49\textwidth]{LebPadova.eps}
    \includegraphics[width=0.49\textwidth]{LebPadova100.eps}
	\caption{Trend of the Lebesgue constant defined in Eq. \eqref{eq:LebConst} for Padua triangles. The number of Friedrichs-Keller triangles is $ N = 2n^2 $, where $n$ is the number of subdivisions of each axis, and the corresponding polynomial degree is $ m $, here shown by red circles.} \label{fig:LebPD}
\end{figure}

\begin{remark}
The theoretical and computational methods developed for triangles in this article can be naturally extended, with small modifications, to other partitionings of the domain $\Omega$, such as %the framework 
in \cite{BruniElefante}. In order to obtain reasonable convergence results, attention has to be paid %mainly 
to the regularity and the uniformity of the elements in the partitions.%, as well as to the convexity of the cells in order to be able to apply the mean value theorem. 

Moreover, the approaches presented in this work can be readily generalized also to higher dimensions using, for instance, simplices instead of triangles. Additionally, beyond the use of Padua points, other explicit point sets with slowly increasing Lebesgue constants such as the Xu points \cite{Xu1996}, the Lissajous node points \cite{Dencker2017,Dencker2017b,Erb2016Degenerate,Erb2016NonDegenerate}, or similar node sets as in \cite{Harris2013} are equally suitable. In the next section, we will further explore alternative strategies based on the Fekete and the Leja triangles.
\end{remark}

\section{Fekete triangles, Leja sequences and their implementation}
\label{sec3new}
In the nodal framework, a common strategy for constructing sets with a low Lebesgue constant consists in looking for Fekete nodes or Leja sequences of nodes. While the theoretical determination of such sets is rather tough \cite{BruniErbFekete}, some linear algebra based algorithms for their approximate computation have been studied \cite{bos2010computing}. We extend these greedy methods to our histopolation setting, and note that, in contrast to Padua triangles, Leja and approximate Fekete triangles may be extracted also from very irregular meshes.

\subsection{Fekete and approximate Fekete triangles}
Fekete nodes are points that maximize the determinant of the Vandermonde matrix. In the univariate case on the interval $[-1,1]$, they coincide with the Legendre-Gauss-Lobatto nodes, and their Lebesgue constant exhibits logarithmic growth \cite{IbrahimogluSurvey}. In higher-dimensional settings, Fekete points maintain a small growth of the Lebesgue constant as the number of interpolation points increases. However, explicit formulas for these points are available only in a few specific cases (see, for instance, \cite{Bos2023}) and one has to stick to numerical approximations, as discussed in \cite{BSV12}. For histopolation, Fekete segments within the interval $[-1,1]$ have been introduced and analyzed in \cite{BruniErbFekete}. Also in this setting, the Lebesgue constant for the Fekete segments grows at a very slow rate. 

Similar as for Fekete segments in $[-1,1]$, one can define a set of Fekete triangles $ \mathcal{T}_M^{\mathrm{Fek}} $ in $\Omega$ as the set of triangles $$\mathcal{T}_M^{\mathrm{Fek}} =\left\{ t_1^{\mathrm{Fek}}, \ldots, t_M^{\mathrm{Fek}} \right\}$$ 
such that the Vandermonde determinant
\begin{equation} \label{eq:FeketeDet}
	\det \left[ V_{i,j} \right] = \det \left[ \frac{1}{|t_i|} \int_{t_i} p_j(x,y) \de x \de y \right], \quad i,j \in \{1, \ldots, M\}, 
	\end{equation}
is maximized over all possible sets $\{t_1, \ldots, t_M\}$ of triangles in $\Omega$.

Similar as in the segmental case \cite{BruniErbFekete}, one can generally show that the corresponding Lebesgue constant grows, at most, as the cardinality of the space $ \P_m\left(\mathbb{R}^2\right)  $:
\begin{equation} \label{eq:LebFek}
	\mathscr{L}_{m} \left(\mathcal{T}_M^{\mathrm{Fek}}\right) \leq \dim\left( \P_m \left(\mathbb{R}^2\right)\right) .
\end{equation}
This inequality arises from basic linear algebra principles, and remains independent of whether a linear functional represents a point evaluation or an average over a domain. Moreover, it does not depend on the domains used for averaging the function. 

To make the maximization of the Vandermonde determinant \eqref{eq:FeketeDet} computationally feasible, we will, as in the case of the Padua triangles, extract approximate Fekete triangles from an initial triangulation $\T_N$. To further reduce the computational cost, we will, as described in \cite{BSV12}, compute these approximate Fekete triangles $\mathcal{T}_M^{\mathrm{Fek}}$ from the triangulation $\T_N$ by applying a column-pivoted QR decomposition to the Vandermonde matrix, which corresponds to an iterative greedy algorithm for the selection of the triangles. The respective algorithm for the numerical extraction of the approximate Fekete triangles, which is almost identical to the nodal one, is given in Algorithm \ref{algFek}. In Figure \ref{LejaeFekTriangle} (left) approximate Fekete triangles extracted from a regular Friedrichs-Keller triangulation are displayed. Note that as the approximate Fekete triangles $\mathcal{T}_M^{\mathrm{Fek}}$ are extracted from a finite collection $\mathcal{T}_N$ the estimate \eqref{eq:LebFek} for the Lebesgue constant may in fact not hold true anymore. This gets visible also in the numerical evaluations of the Lebesgue constants given in Figure \ref{fig:LebLF} and Figure \ref{fig:LebLFdetail}. 

\begin{algorithm}[h!]
	\caption{Extraction of approximate Fekete triangles}
	\begin{algorithmic}[1]
		\Require a triangulation $\mathcal{T}_N= \{t_1,\dots,t_N\}$ of $ \Omega $, a basis $ \left\{p_1, \dots,p_M \right\}$ for $ \P_m \left(\mathbb{R}^2\right) $
		\State Compute the rectangular collocation matrix $ W \in \mathbb{R}^{N \times M}$ given as
        \[ [W_{i,j}] = \left[ \mu_i(p_j)\right] = \left[ \frac{1}{|t_i|} \int_{t_i} p_j(x,y) \de x \de y \right].\]
		\State Calculate the column pivoted QR decomposition of $W^T$ as 
        \[Q R = W^T P, \]
        with a permutation matrix $P \in \mathbb{R}^{N \times N}$. 
        \State Store the first $M$ indices $I = \{i_1, \ldots, i_M\}$ linked to the permutation $P$ of the columns of the matrix $W^T$. 
        \Ensure the set of approximate Fekete triangles $\mathcal{T}^{\mathrm{Fek}}_M= \left\{ t_{i_1},\dots,t_{i_M} \right\}$. 
	\end{algorithmic}
	\label{algFek}
\end{algorithm}

\subsection{Discrete Leja sequences of triangles}
 Instead of performing a global optimization for approximate Fekete triangles based on the Vandermonde determinant in \eqref{eq:FeketeDet}, one can adopt an alternative approach by ordering the basis $ \left\{ p_1, \ldots, p_M \right\} $ of the polynomial space $ \P_m\left(\mathbb{R}^2\right) $ and constructing a sequence $ \mathcal{T}_M^{\mathrm{Leja}} $ of triangles $ \left\{t_1^{\mathrm{Leja}}, \ldots, t_M^{\mathrm{Leja}}\right\} $. At each iteration step $ k \in \{1, \ldots, M\} $, this methods selects a new triangle that maximizes the determinant of a progressively increasing Vandermonde matrix, following the idea of Leja sequences, as for instance described in \cite{DeMarchiLeja}. 
 
 The greedy algorithm for extracting a discrete Leja sequence from the triangulation $\T_N$, similar as in the nodal case, is presented in Algorithm \ref{algLeja}. Figure \ref{LejaeFekTriangle} (right) shows a set of Leja triangles extracted from a regular Friedrichs-Keller triangulation.

\begin{algorithm}[h!]
	\caption{Extraction of discrete Leja triangles}
		\begin{algorithmic}[1]
		\Require a triangulation $\mathcal{T}_N= \{t_1,\dots,t_N\}$ of $ \Omega $, a basis $ \left\{p_1, \dots,p_M \right\}$ for $ \P_m \left(\mathbb{R}^2\right) $
		\State Compute the rectangular collocation matrix $ W \in \mathbb{R}^{N \times M}$ given as
        \[ [W_{i,j}] = \left[ \mu_i(p_j)\right] = \left[ \frac{1}{|t_i|} \int_{t_i} p_j(x,y) \de x \de y \right].\]
		\State Calculate the LU decomposition of $W$ with partial pivoting such that 
        \[L R  = P W, \]
        with a permutation matrix $P \in \mathbb{R}^{N \times N}$. 
        \State Extract the first $M$ indices $I = \{i_1, \ldots, i_M\}$ linked to the permutation $P$ of the rows of the matrix $W$. 
        \Ensure the set of discrete Leja triangles $\mathcal{T}^{\mathrm{Leja}}_M= \left\{ t_{i_1},\dots,t_{i_M} \right\}$. 
	\end{algorithmic}
	\label{algLeja}
\end{algorithm}

\begin{figure}[ht]
	\centering
	\includegraphics[width=0.49\textwidth]{FeketeTriangle2020.eps}
	\includegraphics[width=0.49\textwidth]{LejaTriangle2020.eps}
	\caption{Approximate Fekete (left) and discrete Leja triangles (right) extracted from the regular Friedrichs-Keller triangulation $\mathcal{T}^{\mathrm{reg}}_{800}$ using $n = 20$ subintervals on each axis.}
	\label{LejaeFekTriangle}
\end{figure}

% 
%%and the set of Leja triangles
%\begin{equation*}
%	\mathcal{T}_m^{\mathrm{Leja}}=\left\{t^{\mathrm{Leja}}_{\iota}\, :\,   \iota=1,\dots,\kappa(m)\right\}\subset \mathcal{T}_{N}, 
%\end{equation*}
%by using the Greedy Algorithms~\ref{algFek} and~\ref{algLeja}, respectively~\cite{bos2010computing}. See Figure~\ref{LejaeFekTriangle}.

%{\color{blue} Non mi è chiarissimo come si relaziona il numero di triangoli di Leja/Fekete al grado polinomiale. Devo rifletterci un attimo. L.}

In Figure \ref{fig:LebLF}, the behavior of the Lebesgue constant \eqref{eq:LebConst} for approximate Fekete triangles and discrete Leja sequences of triangles is visualized. The respective computations are based on Vandermonde matrices using tensor-product Chebyshev polynomials as a basis. The selection of Fekete and Leja triangles requires a predetermined choice of the total polynomial degree $ m $. To be consistent with the choice of the Padua triangles, we set $ m $ to be the largest admissible value for the Padua triangles as established in Proposition \ref{thmimp}. 

However, for the Fekete and Leja triangles, this choice of $m$ may sometimes be too large relative to the overall size of the triangulation $\T_N$, leading to a clear violation of the estimate in Eq. \eqref{eq:LebFek}. This phenomenon is evident in Figure \ref{fig:LebLF} for certain isolated values $m$, and becomes even more pronounced at a finer discretization level, as shown in Figure \ref{fig:LebLFdetail}. Despite these specific outliers, the approximate Fekete triangles yield a slightly better Lebesgue constant than the discrete Leja triangles. However, when compared to the Padua triangles (cf. Figure \ref{fig:LebPD}), both Fekete and Leja triangles seem to exhibit an inferior performance. 

\begin{figure}[ht]
	\centering
	\includegraphics[width=0.49\textwidth]{LebFek100.eps}
	\includegraphics[width=0.49\textwidth]{LebLeja100.eps}
	\caption{Lebesgue constants of approximate Fekete (left) and discrete Leja triangles (right) computed on the regular triangulation $\mathcal{T}^{\mathrm{reg}}_{N}$, where $ N = 2 n^2$ and $ n \in \{ 10, 20, \ldots, 100 \}$ is the number of subdivisions of each axis.}
	\label{fig:LebLF}
\end{figure}

\begin{figure}[ht]
	\centering
	\includegraphics[width=0.49\textwidth]{LebFek.eps}
	\includegraphics[width=0.49\textwidth]{LebLeja.eps}
	\caption{Lebesgue constants of Fekete (left) and Leja triangles (right) computed on the regular triangulation $\mathcal{T}^{\mathrm{reg}}_{N}$. Here, all subdivision numbers $n \in\{5, 6, \ldots, 40 \}$ are considered. }
	\label{fig:LebLFdetail}
\end{figure}

\subsection{Calculation of histopolation and histopolation-regression operators}\label{sec3}

In analogy to the histopolation operator \eqref{ea:histopolationPadua} and the histopolation-regression operator \eqref{operatorPhat} for the Padua triangles, we can define corresponding projection operators using the approximate Fekete triangles $\T^{\mathrm{Fek}}_{M}$ and the discrete Leja triangles $\T^{\mathrm{Leja}}_{M}$ instead. While the Padua triangles require as reference domain the square $\Omega = [-1,1]^2$, the Fekete and Leja triangles can be extracted from triangulations $\T_N$ of a general polygonal domain. 

In case of the approximate Fekete triangles $\mathcal{T}_M^{\mathrm{Fek}}$, fixing a basis $\{p_1, \ldots p_M\}$ of $\mathbb{P}_m\left(\mathbb{R}^2\right)$, the histopolation operator ${\Pi}^{\mathrm{Fek}}_{m}$ is given as 
\begin{equation*}\label{operatorinterpFek}
\begin{array}{rcl}
{\Pi}^{\mathrm{Fek}}_{m}: f \in C\left(\Omega\right) &\rightarrow& {\Pi}^{\mathrm{Fek}}_{m}[f] =\sum\limits_{j=1}^{M}a_j p_j\in \mathbb{P}_m\left(\mathbb{R}^2\right),
\end{array}
\end{equation*}
where the vector of coefficients $\left[a_1,\dots,a_{M}\right]^T$ are determined such that the histopolation conditions $\mu_i(f) = \mu_i\left(\sum_{j=1}^{M}a_j p_j \right)$, $i \in \{1, \ldots, M\}$, are satisfied on the Fekete triangles $\mathcal{T}_M^{\mathrm{Fek}}$. On the other hand, for the histopolation-regression problem, we get the projection operator
\begin{equation}\label{operatorPhatFek}
\begin{array}{rcl}
\hat{\Pi}^{\mathrm{Fek}}_{d}: f \in C\left(\Omega\right) &\rightarrow& \hat{\Pi}^{\mathrm{Fek}}_{d}[f] = \sum\limits_{j=1}^{D}\hat{a}_j p_j \in \mathbb{P}_d\left(\mathbb{R}^2\right). 
\end{array}
\end{equation}
where $d > m$ and $\hat{\boldsymbol{a}}=\left[\hat{a}_1,\dots,\hat{a}_{D}\right]^T$ is the solution of the augmented linear system~\eqref{LagrangeMultmethod} relative to the solution of the constrained least-squares problem involving exact histopolation on $\mathcal{T}_M^{\mathrm{Fek}}$ and regression on the remaining triangles. In this case, to set up the system of equations, the triangulation $\T_N$ gets reordered such that $\mathcal{T}_M^{\mathrm{Fek}}$ denote the first $M$ triangles in $\T_N$. 
 In a completely analogous way, the polynomial histopolation operator and the histopolation-regression operator on the Leja triangles $\mathcal{T}_M^{\mathrm{Leja}}$ can be introduced using the notation
\begin{equation*}\label{operatorinterpLeja}
\begin{array}{rcl}
{\Pi}^{\mathrm{Leja}}_{m}: f \in C\left(\Omega\right) &\rightarrow& {\Pi}^{\mathrm{Leja}}_{m}[f] = \sum\limits_{j=1}^{M}a_j p_j\in \mathbb{P}_m\left(\mathbb{R}^2\right),
\end{array}
\end{equation*}
and
\begin{equation}\label{operatorPhatLeja}
\begin{array}{rcl}
\hat{\Pi}^{\mathrm{Leja}}_{d}: f\in C\left(\Omega\right) &\rightarrow& \hat{\Pi}^{\mathrm{Leja}}_{d}[f]=\sum\limits_{j=1}^{D} \hat{a}_j p_j \in \mathbb{P}_d\left(\mathbb{R}^2\right). 
\end{array}
\end{equation}

We notice that the properties shown for the Padua histopolation-regression operator~\eqref{operatorPhat}, including the bound for the Lebesgue constant, can be adapted also for the operators~\eqref{operatorPhatFek} and~\eqref{operatorPhatLeja}. On the other hand, the quasi-optimal logarithmic bound of the Lebesgue constant for histopolation on the Padua triangles seen in Proposition \ref{prop:stability2} seems not to be true for the approximate Fekete and discrete Leja triangles (cf. Figure \ref{fig:LebLF} and Figure \ref{fig:LebLFdetail}).  

We also note that while for the Padua triangles Proposition \ref{prop:stability} gives a theoretic guarantee for the unisolvence of the histopolation problem, the unisolvence in case of the Fekete and Leja triangles has to be checked numerically. This can be done implicitly within the pivoted QR or LU factorization in the extraction of the respective triangles. 

%We introduce two histopolation-regression methods using the same idea of the operator~\eqref{operatorPhat}, which may also be readily adapted to approximate functions defined on any bivariate domain $\Omega$. In particular, let $\Omega \subset \mathbb{R}^2$ be a domain and let $\mathcal{T}_N$, $N\in\mathbb{N}$ be a triangulation of $\Omega$ formed by $N$ triangles. We then consider an unisolvent subset $ \mathcal{T}_M \subset \mathcal{T}_N $ for histopolation in $ \P_m \left(\R^2\right) $, and an intermediate space $ \P_d \left(\R^2\right) $ such that $ D = \dim\left(\P_d \left(\R^2\right)\right)$ and $ M \leq D \leq N $, which we will use for regression purposes. The relationship between $ M $ and $ D $ is studied in the literature, and we will recall it in Section \ref{sec4}.
%In analogy with the definition of the operator~\eqref{operatorPhat}, we set a value $m\in\mathbb{N}$ such that
%\begin{equation}\label{condM}
%    M=\binom{m+2}{2}<N.
%\end{equation}
%Then, we aim to identify an optimal subset of $\mathcal{T}_N$ for the interpolation formed by $\kappa(m)$ triangles and use the remaining triangles to improve the accuracy of the approximation through a simultaneous regression. To this
%aim, we follow the same idea of the \emph{Fekete nodes} and \emph{discrete Leja nodes}. In the classical polynomial interpolation these sets are good sets for interpolation with a low Lebesgue constants.
%Let us describe the former set. Fixed a basis for the polynomial space $\mathbb{P}_d(\mathbb{R})$, Fekete nodes maximize the absolute value of the determinant of the Vandermonde matrix. It is not difficult to show that, in the continuous case, such a collection does not depend on the basis chosen; it is also important to point out that in the discrete case does. To extend this technique to histopolation, it was noted in~\cite{BruniErbFekete} that the normalization in the degrees of freedom~\eqref{eq:dofhistopolation} is necessary. Following such a work, we define the \emph{Fekete triangles} as those supports that maximize
%\begin{equation} \label{eq:VandermondeFekete}
%    DEF .
%\end{equation}
%
%\begin{proposition}
%    For supports $ \{ t_1, \ldots, t_N \} $ that satisfy \eqref{eq:VandermondeFekete}, one has
%    $$ \Lambda_d (t_1, \ldots, t_N) \leq \dim \mathbb{P}_d\left(\mathbb{R}^2\right) .$$
%\end{proposition}
%Already in the nodal case, the identification of such supports is a rather hard problem. Nevertheless, it has been estimated that the corresponding Lebesgue constant is in fact smaller than that predicted by the above inequality \cite{BSV12}. It is also worth noting that, when passing to the discrete problem, such an inequality may be violated. Numerical experiments in~\cite{BruniElefante} showed, in fair generality, that this is not the case. 
%
%
%
%More precisely, by fixing $m$ satisfying~\eqref{condM}, we set $r\in \mathbb{N}$ such that $r>m$. We set
%\begin{equation}
%    \mathcal{B}_r=\left\{e_1(\boldsymbol{x}),\dots,e_{\kappa(r)}(\boldsymbol{x})\right\}, \qquad \boldsymbol{x}=(x,y)\in\Omega, \qquad \kappa(r):=\binom{r+2}{2},
%\end{equation} 
%a basis of the polynomial space $\mathbb{P}_r\left(\mathbb{R}^2\right)$ satisfying the property~\eqref{propimpbas}. 
%In order to define the set of Fekete triangles and the set of Leja triangles, we consider
%\begin{equation*}
%  A=\begin{bmatrix}
%\nu_1\left(e_1\right) & \nu_1\left(e_2\right) & \cdots & \nu_{1}\left(e_{\kappa(m)}\right)\\
%\nu_2\left(e_1\right) & \nu_2\left(e_2\right) & \cdots & \nu_{2}\left(e_{\kappa(m)}\right)\\
%\vdots  & \vdots  & \ddots & \vdots  \\
%\nu_{N}\left(e_1\right) & \nu_{N}\left(e_2\right) & \cdots & \nu_{N}\left(e_{\kappa(m)}\right)\\
%\end{bmatrix}.
%\end{equation*}
%We compute the set of Fekete triangles
%\begin{equation*}
%\mathcal{T}_{\kappa(m)}^{\mathrm{Fek}}=\left\{t^{\mathrm{Fek}}_{\iota}\, :\,   \iota=1,\dots,\kappa(m)\right\}\subset \mathcal{T}_{N}, 
%\end{equation*}
%and the set of Leja triangles
%\begin{equation*}
%\mathcal{T}_{\kappa(m)}^{\mathrm{Leja}}=\left\{t^{\mathrm{Leja}}_{\iota}\, :\,   \iota=1,\dots,\kappa(m)\right\}\subset \mathcal{T}_{N}, 
%\end{equation*}
% by using the Greedy Algorithms~\ref{algFek} and~\ref{algLeja}, respectively~\cite{bos2010computing}. See Figure~\ref{LejaeFekTriangle}.
%\begin{algorithm}
%\caption{}
%\begin{algorithmic}[]
%\Require $\mathcal{T}_N=\left[t_1,\dots,t_N\right],  \mathcal{B}_r=\left\{\boldsymbol{e}_{0},\dots,\boldsymbol{e}_{\kappa(m)}\right\}$ 
%\Ensure $\mathcal{T}^{\mathrm{Fek}}_m=\left[t_1^{\mathrm{Fek}},\dots,t_{\kappa(m)}^{\mathrm{Fek}}\right]$
%\State Compute $N$ and $\kappa(m)$
%\State $\mathrm{A}_0=A$; $k=1:N$
%\State $[Q_0,R_0,\Sigma_0]=\mathrm{qr}\left(\mathrm{A}^{\prime}_0\right)$; $k=\Sigma_0 k;$
%\State $\mathcal{T}^{\mathrm{Fek}}_m=\left[t_1^{\mathrm{Fek}},\dots,t_{\kappa(m)}^{\mathrm{Fek}}\right]=\mathcal{T}_N(k_1,\dots,k_{\kappa(m)})$.
%\end{algorithmic}
%\label{algFek}
%\end{algorithm}
%
%\begin{algorithm}
%\caption{}
%\begin{algorithmic}[]
%\Require $\mathcal{T}_N=\left[t_1,\dots,t_N\right],  \mathcal{B}_r=\left\{\boldsymbol{e}_{0},\dots,\boldsymbol{e}_{\kappa(m)}\right\}$ 
%\Ensure $\mathcal{T}^{\mathrm{Leja}}_m=\left[t_1^{\mathrm{Leja}},\dots,t_{\kappa(m)}^{\mathrm{Leja}}\right]$
%\State Compute $N$ and $\kappa(m)$
%\State $\mathrm{A}_0=A$; $k=1:N$
%\State $[L_0,U_0,\Sigma_0]=\mathrm{lu}\left(\mathrm{A}_0\right)$; $k=\Sigma_0 k;$
%\State $L_m=\left[t_1^{\mathrm{Leja}},\dots,t_{\kappa(m)}^{\mathrm{Leja}}\right]=\mathcal{T}_N(k_1,\dots,k_{\kappa(m)})$.
%\end{algorithmic}
%\label{algLeja}
%\end{algorithm}
%By following the same line of the previous section, we assume that the triangulation  $\mathcal{T}_{N}$ is reordered so that 
%\begin{equation}\label{reordFek}
%t_{i}=t^{\mathrm{Fek}}_{i}, \qquad i=1,\dots,M.
%\end{equation}
%Under this assumption, by reordering the histopolating linear functionals accordingly, 
%we define the polynomial histopolation operator on the Fekete triangles $\mathcal{T}_M^{\mathrm{Fek}}$ as
%\begin{equation*}\label{operatorinterpFek} \begin{array}{rcl}{\Pi}^{\mathrm{Fek}}_{m}: f(\boldsymbol{x})\in C\left(\Omega\right) &\rightarrow& {\Pi}^{\mathrm{Fek}}_{m}[f](\boldsymbol{x})=\sum\limits_{i=1}^{M}a_ie_i(\boldsymbol{x})\in \mathbb{P}_m\left(\mathbb{R}^2\right), \qquad \boldsymbol{x}\in \Omega, 
%\end{array}
%\end{equation*}
%where the vector of coefficients $\left[a_1,\dots,a_{M}\right]^T$ is computed such that 
%\begin{equation*}
%\int_{t_i} f(x,y) \de x \de y  = \int_{t_i} p(x,y) \de x \de y, \qquad i=1,\dots,M.
%\end{equation*}
%Finally, we define the Fekete histopolation-regression operator as
%\begin{equation}\label{operatorPhatFek}
%\begin{array}{rcl}
%\hat{\Pi}^{\mathrm{Fek}}_{d}: f(\boldsymbol{x})\in C\left(\Omega\right) &\rightarrow& \hat{\Pi}^{\mathrm{Fek}}_{d}[f](\boldsymbol{x})=\sum\limits_{i=1}^{D}\hat{a}_ie_i(\boldsymbol{x})\in \mathbb{P}_d\left(\mathbb{R}^2\right), \qquad \boldsymbol{x}\in \Omega, 
%\end{array}
%\end{equation}
%where $\hat{\boldsymbol{a}}=\left[\hat{a}_1,\dots,\hat{a}_{D}\right]^T$ is the solution of the KKT linear equations~\eqref{LagrangeMultmethod}. In a manner completely analogous to previous methods, we assume a reordering of the triangulation $\mathcal{T}_N$ similar to~\eqref{reordFek}, followed by a reordering of the corresponding linear functionals. We then define the polynomial histopolation operator on the Leja triangles $\mathcal{T}_M^{\mathrm{Leja}}$ as
%\begin{equation*}\label{operatorinterpLeja}
%\begin{array}{rcl}
%{\Pi}^{\mathrm{Leja}}_{m}: f(\boldsymbol{x})\in C\left(\Omega\right) &\rightarrow& {\Pi}^{\mathrm{Leja}}_{m}[f](\boldsymbol{x})=\sum\limits_{i=1}^{M}a_ie_i(\boldsymbol{x})\in \mathbb{P}_m\left(\mathbb{R}^2\right), \qquad \boldsymbol{x}\in \Omega.
%\end{array}
%\end{equation*}
%Under this assumption, in analogy to~\eqref{operatorPhat} and~\eqref{operatorPhatFek}, we can define the Leja histopolation-regression operator as
%\begin{equation}\label{operatorPhatLeja}
%\begin{array}{rcl}
%\hat{\Pi}^{\mathrm{Leja}}_{d}: f(\boldsymbol{x})\in C\left(\Omega\right) &\rightarrow& \hat{\Pi}^{\mathrm{Leja}}_{d}[f](\boldsymbol{x})=\sum\limits_{i=1}^{D}\hat{a}_ie_i(\boldsymbol{x})\in \mathbb{P}_d\left(\mathbb{R}^2\right), \qquad \boldsymbol{x}\in \Omega.
%\end{array}
%\end{equation}

Algorithm \ref{algfinal} summarizes the calculation of the histopolation-regression polynomial $\hat{\Pi}_d^{\mathrm{Pad}}[f]$ for the Padua triangles. The respective adaption to the operators $\hat{\Pi}^{\mathrm{Fek}}_{d}$ and $\hat{\Pi}^{\mathrm{Leja}}_{d}$ defined in~\eqref{operatorPhatFek} and~\eqref{operatorPhatLeja} is immediate. 

\begin{algorithm}[h!]
	\caption{histopolation-regression operators}
	\begin{algorithmic}[1]
		\Require \begin{itemize}
		    \item[a)] $m\in\mathbb{N}$,  $d\in\mathbb{N}$, a triangulation $\mathcal{T}_N= \{t_1,\dots,t_N\}$ of $ \Omega $,
            \item[b)] a basis $ \left\{p_1, \dots,p_D \right\}$ for $ \P_d \left(\mathbb{R}^2\right) $ with $\left\{p_1,\dots,p_M\right\}$ being a basis for $\mathbb{P}_m\left(\mathbb{R}^2\right)$,
            \item[c)] the data vector $\boldsymbol{b} = [\mu_1(f), \ldots, \mu_N(f)]^T$.
		\end{itemize}

\vspace{2mm}
        \State Compute the Padua triangles $\T_M^{\mathrm{Pad}}$ by Algorithm~\ref{algPad}.
        \State Reorder $\mathcal{T}_N$ such that $\T_M^{\mathrm{Pad}}$ are the first $M$ elements of $\mathcal{T}_N$, reorder respectively $\boldsymbol{b}$ such that the first $M$ entries $\boldsymbol{d}$ contain the data values on the triangles $\T_M^{\mathrm{Pad}}$.  
        \State Compute the collocation matrices $ W $ and $ C $ (using a Gaussian quadrature rule)
		and create the augmented matrix $\begin{bmatrix}
		2 W^T W & C^T  \\
		C & 0  \\
	\end{bmatrix}$ for the linear system~\eqref{LagrangeMultmethod}.
		\State Compute the coefficient vector $\hat{\boldsymbol{a}}=\left[\hat{a}_1,\dots,\hat{a}_D\right]^T$ by solving the linear system~\eqref{LagrangeMultmethod}.
		\State Calculate the polynomial $\sum\limits_{j=1}^D \hat{a}_j p_j$.
        \Ensure the approximation polynomial  $\hat{\Pi}_d^{\mathrm{Pad}}[f] = \sum\limits_{j=1}^D \hat{a}_j p_j$.
	\end{algorithmic}
	\label{algfinal}
\end{algorithm}


%
%\begin{remark}
%   We observe that the sets of Fekete and Leja triangles can be identified from a triangulation of any domain $\Omega$, see Figure~\ref{LejaeFekTriangleontriangle}. In contrast, the set of Padua triangles can only be extracted from a triangulation of a rectangular domain $\Omega=[a,b]\times[c,d]$. Therefore,  operators~\eqref{operatorPhatFek} and~\eqref{operatorPhatLeja} can be used for any domain, whereas operator~\eqref{operatorPhat} is limited to rectangular domains.
%\end{remark}
%
%
%
%\begin{figure}[ht]
%    \centering
%        \includegraphics[scale=0.49]{FeketeTriangle2020.eps}
%    \includegraphics[scale=0.49]{LejaTriangle2020.eps}
%    \caption{Fekete triangles (left) and Leja triangles (right) in red computed relative to the regular triangulation $\mathcal{T}^{\mathrm{reg}}_{20,20}$.}
%    \label{LejaeFekTriangle}
%\end{figure}
%
%
%\begin{figure}[ht]
%    \centering
%        \includegraphics[scale=0.49]{trianglesFekete.eps}
%    \includegraphics[scale=0.49]{trianglesLeja.eps}
%    \caption{Fekete triangles (left) and Leja triangles (right) in red computed relative to the regular triangulation $\mathcal{T}^{\mathrm{reg}}_{20,20}$.}
%\label{LejaeFekTriangleontriangle}
%\end{figure}

%\section{Algorithms}
%\label{secalg}
%In this section, we present the main algorithms used to implement the histopolation-regression operators discussed in this paper. First, we introduce Algorithm~\ref{algPad}, which extracts the Padua triangle $\T_m^{\mathrm{Pad}}$ from the triangulation $\T_N$. % defined in~\eqref{triangulation}. 
%Once the hypothesis of Proposition \ref{prop:stability} on the maximal triangle length $h_{\max}$ is satisfied, the following Algorithm \ref{algPad} is well-posed and extracts the Padua triangles.
%
%
%Then, Algorithm \ref{algFek} and Algorithm \ref{algLeja} generalize the idea of the approximate Fekete points and discrete Leja points extending them to the case of triangles via the following greedy strategies.
%
%
%
%Finally, Algorithm~\ref{algfinal} computes the histopolation-regression operator $\hat{\Pi}_d^{\mathrm{Pad}}$, defined in~\eqref{operatorPhat}. 
%This algorithm can also be adapted to compute the operators $\hat{\Pi}^{\mathrm{Fek}}_{d}$ and $\hat{\Pi}^{\mathrm{Fek}}_{d}$ defined in~\eqref{operatorPhatFek} and~\eqref{operatorPhatLeja}, respectively. 


\section{Numerical results}
\label{sec4}

In the following, we provide a numerical evaluation of the performance of the histopolation and histopolation-regression methods based on the Padua, Fekete and Leja triangles. For this evaluation, we consider the following three test functions:
\begin{equation*}
    f_1(x,y)=e^{x+y} \sin{(\pi x y)},
    \qquad f_2(x,y)=|x+y|,
    \qquad f_3(x,y)=\frac{1}{1+10(x^2+y^2)}.
\end{equation*}
The first function is regular and smooth, while the second is non-differentiable at the diagonal points $x = -y$. In the second case, we expect lower convergence rates for the histopolation polynomials. The third function is the Runge function already encountered in Figure \ref{fig:Runge}. Note that $ \Vert f_1 \Vert_\infty \approx 4.71 $, $ \Vert f_2 \Vert_\infty = 2 $, and $ \Vert f_3 \Vert_\infty = 1 $. All the errors in the numerical experiments are computed with respect to the uniform norm on $[-1,1]^2$.

%In this section, we aim to prove the accuracy of our proposed methods through three types of experiments. In particular, we consider the following test functions
%\begin{equation*}
%    f_1(x,y)=\frac{1}{x+y+5}, \qquad f_2(x,y)=e^{x+y}, \qquad f_3(x,y)=\sin(x+y).
%\end{equation*}

\subsection{Comparison of histopolators}
We compare the results obtained for the Padua triangles with the histopolators based on the approximate Fekete triangles and the discrete Leja triangles. As a rule for determining the degree $m$ of the polynomial space $\P_m\left(\mathbb{R}^2 \right)$, the number $m$ is chosen largest possible such that the well-posedness condition of Proposition \ref{thmimp} is still satisfied. Integrals are computed with a Gaussian quadrature rule. For the test function $ f_1 $, convergence of the schemes can be observed for all three histopolators, as shown in Figure \ref{fig:ConvHist}.

\begin{figure}[H]
    \centering
    \includegraphics[width=0.60\linewidth]{SmoothHistLogy.eps} 
    \caption{Convergence of the histopolators towards the function $ f_1$.}
    \label{fig:ConvHist}
\end{figure}

In contrast, we see that the instability of the Lebesgue constant associated with the Fekete triangles (see Figure \ref{fig:LebLF}) has an impact also for the convergence of the methods, which is not visible anymore for the functions $ f_2 $ and $ f_3 $. This is depicted in Figure \ref{fig:Divergence}, where the error associated with the Fekete histopolator is dashed. A softer instability is shown by the histopolator associated with the Leja triangles. The peaks that can be seen for the Fekete (and partially for the Leja) histopolator correspond to the peaks already observed for the Lebesgue constants in Figure \ref{fig:LebLF} and Figure \ref{fig:LebLFdetail}. 

\begin{figure}[H]
    \centering
    \includegraphics[width=0.49\linewidth]{AbsDashed.eps}
    \
    \includegraphics[width=0.49\linewidth]{Runge10Convergence.eps}
    \caption{Convergence of the histopolators towards the functions $ f_2 $ and $ f_3 $.}% {\color{red}W: ricontrolliamo: per la funzione di Runge mi aspetterei una convergenza piu' visibile, sopratutto per i triangoli di Padova.}}
    \label{fig:Divergence}
\end{figure}

%Results are reported in Figure \ref{fig:ConvHist}.

% \begin{figure}[H]
%     \centering
%     \includegraphics[width=0.49\linewidth]{SmoothHist.eps} \
%     \includegraphics[width=0.49\linewidth]{AbsHist.eps}
%     \\
%     \includegraphics[width=0.49\linewidth]{RungeHist.eps}
%     \caption{Convergence of the histopolators towards the functions $ f_1$, $ f_2$ and $ f_3 $.}
%     \label{fig:ConvHist}
% \end{figure}

As evident from Figure \ref{fig:ConvHist}, the histopolator associated with the Padua triangles is the only one, among those discussed, that shows convergence for all three test functions. %Further, it is also interesting to observe that the histopolator associated with Fekete triangles presents a strong instability when the supports are extracted from a too coarse mesh. This was already observed in the analysis of the associated Lebesgue constant, see Figure \ref{fig:LebLF}.

\subsection{Adding the regression term}

As in the previous experiment, let $m$ be the maximal natural number still satisfying the condition ~\eqref{gencond} of Proposition \ref{thmimp}. Now, we additionally assume for the combined histopolation-regression scheme, similarly as in~\cite{bruni2024polynomial}, that the polynomial degree $d$ satisfies 
\begin{equation*}
    d \sim m+\sqrt{m}.
\end{equation*}
This particular relationship between $m$ and $d$ has its origins in the nodal framework \cite{DeMarchi:2015:OTC}, and we use it for the additional regression term discussed in Section \ref{sect:regressionmethod}. The respective convergence results are presented in Figure \ref{fig:ConvRegr} for $ f_1 $ and $f_2$, and in Figure \ref{fig:ConvRegrRunge} for $ f_3 $.
%To ensure high precision in the integrals, we use a Gaussian quadrature formula with an appropriate number of integration points.

\begin{figure}[H]
    \centering
    \includegraphics[width=0.49\linewidth]{SmoothRegr.eps} \
    \includegraphics[width=0.49\linewidth]{AbsRegr.eps}
    % \\
    % \includegraphics[width=0.49\linewidth]{RungeRegr.eps}
    \caption{Convergence of the histopolation-regression approximants towards the functions $ f_1$ and $ f_2$.}
    \label{fig:ConvRegr}
\end{figure}

We notice that for the smooth function $ f_1 $ the convergence of the histopolants is improved (cf. with Figure \ref{fig:ConvHist}). Furthermore, for the function $ f_2 $, the instabilities and peaks visible in the error of the Fekete histopolation seem to be resolved. A similar effect is visible also for the Runge function $f_3$, see Figure \ref{fig:ConvRegrRunge}. 

%In such a case, however, we do not achieve a monotonic convergence as the grid is refined. Nevertheless, the divergence effect presented by the Fekete histopolator and visible in Figure \ref{fig:Divergence} does not appear anymore.

\begin{figure}[H]
    \centering
    % \includegraphics[width=0.49\linewidth]{SmoothRegr.eps} \
    % \includegraphics[width=0.49\linewidth]{AbsRegr.eps}
    % \\
    \includegraphics[width=0.60\linewidth]{RungeRegr10.eps}
    \caption{Convergence of the histopolation-regression approximants towards the Runge function $ f_3 $.}
    \label{fig:ConvRegrRunge}
\end{figure}

From these numerical tests, we can infer that the regression term seems to accelerate the convergence of the already convergent histopolants, as for instance seen for the smooth function $ f_1 (x,y) $. Also for the functions $ f_2 (x,y) $ and $ f_3 (x,y) $, the additional regression term seems to stabilize and improve the convergence of the approximants associated with the discrete Leja triangles and the approximate Fekete triangles.

\subsection{Random mesh}

As seen in the previous sections, the Padua triangles require a strong hypothesis on the mesh regularity in order to guarantee a well-defined histopolation scheme. In contrast, the triangles given by the greedy approaches in Section \ref{sec3new} can be computed easily and provide well-defined histopolation schemes also for less regular initial triangulations. To see this, we increase in this section the randomness of the initial mesh,  %(in the sense that the grid on the axes is not uniform anymore, but generated randomly) and take the corresponding Delaunay triangulation. We thus 
and then apply Algorithms 1--3 to study the convergence behavior of the resulting histopolation operators.

\subsubsection{Grid from random points on the boundary}

As a first attempt, we consider a non-uniform partitioning of the axes to generate a triangular mesh, as visualized in Figure \ref{fig:MeshRandom}. Observe that very different triangles are selected by the three procedures.

\begin{figure}[H]
    \centering
    % \includegraphics[width=0.49\linewidth]{SmoothRegr.eps} \
    % \includegraphics[width=0.49\linewidth]{AbsRegr.eps}
    % \\
    \includegraphics[width=0.60\linewidth]{RandomMeshSelectionSmall.eps}
    \caption{A randomized Friedrichs-Keller triangulation. The extracted triangles depending on the method are pictured in different colours: blue for Leja, orange for Fekete, yellow for Padua.}
    \label{fig:MeshRandom}
\end{figure}

In this experiment, approximately 30\% of the attempts failed to produce a Padua histopolator, meaning that the attribution map \eqref{eq:PPtoPT} is not well defined. Figure \ref{fig:MeshRandom} shows that the condition on $h_\mathrm{max}$ is frequently violated by this approach. Nevertheless, when produced, the histopolator associated with the Padua triangles performs efficiently. This is shown in Figure \ref{fig:ConvSemiRand}, where a comparison of the histopolation-regression approximants for the three triangle extraction procedures is presented for the functions $ f_1 $ and $ f_3 $.


\begin{figure}[H]
    \centering
    \includegraphics[width=0.49\linewidth]{SmoothSemirand.eps} \
    \includegraphics[width=0.49\linewidth]{RungeSemirand.eps}
    % \\
    % \includegraphics[width=0.49\linewidth]{RungeRegr.eps}
    \caption{Convergence of the histopolation-regression approximants towards the functions $ f_1 $ and $ f_3$, for randomized meshes as given in Figure \ref{fig:MeshRandom}.}
    \label{fig:ConvSemiRand}
\end{figure}

\subsubsection{Random Delaunay triangulation}
A more involved situation emerges when the mesh is completely random, and does not stem from a random subdivision of the axes.
In this case, the attribution map \eqref{eq:PPtoPT} rarely turns out to be well defined, while the extraction of Fekete and Leja triangles is granted. An example of this is given in Figure \ref{fig:MeshRandomComparison}. In this example, the entire triangulation consists of about $2000$ elements, and according to the criteria in Section \ref{sec2}, a regular Friedrichs-Keller triangulation of this size would allow a well-defined histopolation of total degree $ m = 7 $. The dimension of the corresponding space is $ \dim \P_7 (\Omega) = 36 $, but only $ 25 $ Padua triangles (shaded in yellow) are detected. As a consequence, the histopolator is not defined. The right hand panel of that same image shows that $ 36 $ Leja triangles are correctly extracted.

\begin{figure}[H]
    \centering
    \includegraphics[width=0.49\linewidth]{randomPadm7.eps} \
    \includegraphics[width=0.49\linewidth]{randomLejam7.eps}
    \\
%    \includegraphics[width=0.49\linewidth]{RandomRunge.eps}
    \caption{Extraction of Padua and Leja triangles from a random mesh of about 2000 elements.}
    \label{fig:MeshRandomComparison}
\end{figure}

We finally drop the construction of Padua triangles, and compare the convergence of the histopolants associated with the Fekete and the Leja triangles on random meshes. The respective results are quite consistent with the situation pictured for structured mesh. In this case, the randomness of the mesh is partially compensated by the large number of triangles at play. As a consequence, for $ f_1 $ and $ f_2 $ we observe a slow convergence towards the exact functions, see Figure \ref{fig:MeshRandomConvergence}.

\begin{figure}[H]
    \centering
     \includegraphics[width=0.49\linewidth]{SmoothRand.eps} \
     \includegraphics[width=0.49\linewidth]{AbsRand.eps}
    % \\
%    \includegraphics[width=0.49\linewidth]{RandomRunge.eps}
    \caption{Convergence of the histopolators arising from a random mesh as in Figure \ref{fig:MeshRandomComparison}.}
    \label{fig:MeshRandomConvergence}
\end{figure}

% In the first experiment, we examine the trend of the maximum approximation errors
% \begin{eqnarray}    \hat{e}^{\mathrm{Pad}}_d\left(f_i\right)&=&\max_{(x,y)\in \Omega}\left\lvert f_i(x,y)-\hat{\Pi}_{d}^{\mathrm{Pad}}\left[f_i\right](x,y)\right\rvert,\qquad i=1,2,3, \label{err1}\\   \hat{e}^{\mathrm{Fek}}_d\left(f_i\right)&=&\max_{(x,y)\in \Omega}\left\lvert f_i(x,y)-\hat{\Pi}_{d}^{\mathrm{Fek}}\left[f_i\right](x,y)\right\rvert, \qquad i=1,2,3,\label{err2}\\  \hat{e}^{\mathrm{Leja}}_d\left(f_i\right)&=&\max_{(x,y)\in \Omega}\left\lvert f_i(x,y)-\hat{\Pi}_{d}^{\mathrm{Leja}}\left[f_i\right](x,y)\right\rvert, \qquad i=1,2,3, \label{err3}
% \end{eqnarray}
% produced by the operators $\hat{\Pi}_{d}^{\mathrm{Pad}}$, $\hat{\Pi}_{d}^{\mathrm{Fek}}$ and $\hat{\Pi}_{d}^{\mathrm{Leja}}$ for different triangulations. More precisely, we consider $\Omega=[-1,1]^2$ and we consider the following regular triangulations $\mathcal
% {T}^{\mathrm{reg}}_{n,n}$, for different values of $n$, ranging from $n=20$ and $n=100$. The results are shown in Figure~\ref{funf1}. From the plot, we observe that the approximation error decreases as the number of triangles in the triangulation increases, eventually reaching machine precision. Once this accuracy is achieved, the error remains constant even as the number of triangles increases further.
% \begin{figure}[ht]
%     \centering
%         \includegraphics[scale=0.32]{f1.eps}
%     \includegraphics[scale=0.32]{f2.eps}
%     \includegraphics[scale=0.32]{f3.eps}
%     \caption{Semilog plot of the maximum approximation error produced by $\hat{\Pi}_{d}^{\mathrm{Pad}}$, $\hat{\Pi}_{d}^{\mathrm{Fek}}$, and $\hat{\Pi}_{d}^{\mathrm{Leja}}$ for the functions $f_1$ (left), $f_2$ (center), and $f_3$ (right) over the domain $\Omega = [-1,1]^2$ using the regular triangulations $\mathcal{T}^{\mathrm{reg}}_{N}$ with $n = 20, 40, 60, 80, 100$.}
% \label{funf1}
% \end{figure}


% In the second experiment, we compare the errors produced by the histopolation-regression operators~\eqref{operatorPhat},~\eqref{operatorPhatFek} and~\eqref{operatorPhatLeja} 
% with those obtained by the polynomial histopolation on the Padua, Fekete, and Leja triangles, respectively. Specifically, we compare the errors~\eqref{err1},~\eqref{err2} and~\eqref{err3}, with
% \begin{eqnarray}
%     {e}^{\mathrm{Pad}}_m\left(f_i\right)&=&\max_{(x,y)\in \Omega}\left\lvert f_i(x,y)-{\Pi}_{m}^{\mathrm{Pad}}\left[f_i\right](x,y)\right\rvert,\qquad i=1,2,3, \label{err1p}\\   {e}^{\mathrm{Fek}}_m\left(f_i\right)&=&\max_{(x,y)\in \Omega}\left\lvert f_i(x,y)-{\Pi}_{m}^{\mathrm{Fek}}\left[f_i\right](x,y)\right\rvert, \qquad i=1,2,3,\label{err2p}\\  {e}^{\mathrm{Leja}}_m\left(f_i\right)&=&\max_{(x,y)\in \Omega}\left\lvert f_i(x,y)-{\Pi}_{m}^{\mathrm{Leja}}\left[f_i\right](x,y)\right\rvert, \qquad i=1,2,3, \label{err3p}
% \end{eqnarray}
% where $\Pi_m^{\mathrm{Pad}}\left[f_i\right]$, $\Pi_m^{\mathrm{Fek}}\left[f_i\right]$, and $\Pi_m^{\mathrm{Leja}}\left[f_i\right]$ are the histopolation polynomials for $f_i$ on the Padua, Fekete, and Leja triangles, respectively. These comparisons were conducted over the domain $\Omega=[-1,1]^2$ using three different Delaunay triangulations constituting of $N=306$, $N=2650$ and $N=23576$  tringles with no angle smaller than $20^{\circ}$, see Figure~\ref{funf1tr}. The results are shown in Tables~\ref{tab:t1},~\ref{tab:t2} and~\ref{tab:t3}. From these tables, we observe that leveraging the remaining triangles of the triangulation through a simultaneous regression consistently improves the accuracy of the approximation, sometimes significantly, even as the degree of the histopolation polynomial increases.



% \begin{figure}[ht]
%     \centering
%         \includegraphics[scale=0.32]{padscatt.eps}
%             \includegraphics[scale=0.32]{Feketetr.eps}
%     \includegraphics[scale=0.32]{lejatr.eps}
%     \caption{Plot of a Delaunay triangulation consisting of $N=306$ triangles over the domain $\Omega = [-1,1]^2$. The red triangles represent the corresponding Padua triangles (left), Fekete triangles (center), and Leja triangles (right).}
% \label{funf1tr}
% \end{figure}


% \begin{table}
%     \centering
%     \begin{tabular}{c  c  c }
%   & ${e}^{\mathrm{Pad}}_m$ &   $\hat{e}^{\mathrm{Pad}}_r$   \\ \hline
%         $f_1(x)$ &  3.9e-05 &  3.4e-06\\ \hline
%        $f_2(x)$ &  1.6e-03 &   4.5e-05 \\ \hline
%         $f_3(x)$ &  1.4e-03 &  4.1e-05\\ \hline
%     \end{tabular}
%         \begin{tabular}{c  c  c }
%   & ${e}^{\mathrm{Fek}}_m$ &   $\hat{e}^{\mathrm{Fek}}_r$  \\ \hline
%         $f_1(x)$ &  2.8e-05 &  3.2e-06\\ \hline
%        $f_2(x)$ &  2.4e-03 &  5.7e-05 \\ \hline
%         $f_3(x)$ &  1.6e-03 &  4.1e-05\\ \hline
%     \end{tabular}
%     \begin{tabular}{c  c  c }
%   & ${e}^{\mathrm{Leja}}_m$ &   $\hat{e}^{\mathrm{Leja}}_r$   \\ \hline
%         $f_1(x)$ &  4.6e-05 &  4.1e-06\\ \hline
%        $f_2(x)$ &  1.9e-03 &   6.0e-05 \\ \hline
%         $f_3(x)$ &  1.6e-03 &  4.9e-05\\ \hline
%     \end{tabular}
%     \caption{Comparisons between the errors $\hat{e}_r^{\mathrm{Pad}}$, $\hat{e}_r^{\mathrm{Fek}}$ and $\hat{e}_r^{\mathrm{Leja}}$ with ${e}_m^{\mathrm{Pad}}$, ${e}_m^{\mathrm{Fek}}$ and ${e}_m^{\mathrm{Leja}}$ computed over the domain $\Omega=[-1,1]^2$ and a Delaunay triangulations of $N=306$ tringles with no angle smaller than $20^{\circ}$.
% }
%     \label{tab:t1}
% \end{table}



% \begin{table}
%     \centering
%     \begin{tabular}{c  c  c }
%   & ${e}^{\mathrm{Pad}}_m$  &   $\hat{e}^{\mathrm{Pad}}_r$    \\ \hline
%         $f_1(x)$ &  2.4e-07 &  5.5e-09\\ \hline
%        $f_2(x)$ &  1.5e-06 &   2.2e-09 \\ \hline
%         $f_3(x)$ &  1.3e-07 &  2.1e-09\\ \hline
%     \end{tabular}
%         \begin{tabular}{c  c  c }
%   & ${e}^{\mathrm{Fek}}_m$  &   $\hat{e}^{\mathrm{Fek}}_r$   \\ \hline
%         $f_1(x)$ &  2.3e-07 &  7.1e-09\\ \hline
%        $f_2(x)$ &  3.2e-06 &  4.0e-09 \\ \hline
%         $f_3(x)$ &  5.3e-07 &  3.2e-09\\ \hline
%     \end{tabular}
%     \begin{tabular}{c  c  c }
%   & ${e}^{\mathrm{Leja}}_m$  &   $\hat{e}^{\mathrm{Leja}}_r$    \\ \hline
%         $f_1(x)$ &  3.2e-07 &  7.4e-09\\ \hline
%        $f_2(x)$ &  2.1e-06 &   3.6e-09 \\ \hline
%         $f_3(x)$ &  2.4e-07 &  2.9e-09\\ \hline
%     \end{tabular}
%     \caption{Comparisons between the errors $\hat{e}_r^{\mathrm{Pad}}$, $\hat{e}_r^{\mathrm{Fek}}$ and $\hat{e}_r^{\mathrm{Leja}}$ with ${e}_m^{\mathrm{Pad}}$, ${e}_m^{\mathrm{Fek}}$ and ${e}_m^{\mathrm{Leja}}$ computed over the domain $\Omega=[-1,1]^2$ and a Delaunay triangulations of $N=2650$ tringles with no angle smaller than $20^{\circ}$. 
% }
%     \label{tab:t2}
% \end{table}


% \begin{table}
%     \centering
%     \begin{tabular}{c  c  c }
%   & ${e}^{\mathrm{Pad}}_m$ &   $\hat{e}^{\mathrm{Pad}}_r$   \\ \hline
%         $f_1(x)$ &  9.7e-10 &  2.4e-11\\ \hline
%        $f_2(x)$ &  1.4e-10 &   7.0e-14 \\ \hline
%         $f_3(x)$ &  8.4e-12 &  5.9e-14\\ \hline
%     \end{tabular}
%         \begin{tabular}{c  c  c }
%   & ${e}^{\mathrm{Fek}}_m$ &   $\hat{e}^{\mathrm{Fek}}_r$  \\ \hline
%         $f_1(x)$ &  7.4e-10 &  7.9e-12\\ \hline
%        $f_2(x)$ &  3.1e-10 &  8.6e-14 \\ \hline
%         $f_3(x)$ &  8.1e-12 &  3.7e-14\\ \hline
%     \end{tabular}
%     \begin{tabular}{c  c  c }
%   & ${e}^{\mathrm{Leja}}_m$ &   $\hat{e}^{\mathrm{Leja}}_r$   \\ \hline
%         $f_1(x)$ &  1.9e-09 &  8.9e-12\\ \hline
%        $f_2(x)$ &  3.1e-10 &   1.2e-13 \\ \hline
%         $f_3(x)$ &  1.1e-11 &  1.3e-13\\ \hline
%     \end{tabular}
%     \caption{Comparisons between the errors $\hat{e}_r^{\mathrm{Pad}}$, $\hat{e}_r^{\mathrm{Fek}}$ and $\hat{e}_r^{\mathrm{Leja}}$ with ${e}_m^{\mathrm{Pad}}$, ${e}_m^{\mathrm{Fek}}$ and ${e}_m^{\mathrm{Leja}}$ computed over the domain $\Omega=[-1,1]^2$ and a Delaunay triangulations of $N=23576$ tringles with no angle smaller than $20^{\circ}$.  
% }
%     \label{tab:t3}
% \end{table}

% In all experiments discussed above, we used the highest value of $m$ satisfying~\eqref{gencond} of Theorem~\ref{thmimp}. In the final experiment, we show that this value provides the best approximation compared to a smaller $m$. Since we observe from both types of experiments that the approximations produced by our three methods are comparable, yielding similar errors and achieving approximately the same accuracy, we focus on the approximation produced by the first method. Specifically, we analyze the behavior of the approximation error produced by $\Pi_m^{\mathrm{Pad}}$ over the domain $\Omega = [-1,1]^2$ for various values of $m$ satisfying~\eqref{gencond} of Theorem~\ref{thmimp}, using the regular triangulation $\mathcal{T}_{n,n}^{\mathrm{reg}}$ with $n = 20, 40, 80$.  The results of these tests are shown in Figure~\ref{funf123n}. From the plots, we observe that the error consistently decreases as the degree $m$ increases. Therefore, it is advantageous to use the highest possible value of $m$ that satisfies Theorem~\ref{thmimp}.

% \begin{figure}[ht]
%     \centering
%         \includegraphics[scale=0.32]{f123n20.eps}
%     \includegraphics[scale=0.32]{f123n40.eps}
%     \includegraphics[scale=0.32]{f123n80.eps}
%     \caption{Semilog plot of the maximum approximation error produced by $\Pi^{\mathrm{Pad}}_m$ over the domain $\Omega=[-1,1]^2$, for any $m$ satisfying~\eqref{conditionOnm} of Theorem~\ref{thmimp}, using the regular triangulation $\mathcal{T}_{n,n}^{\mathrm{reg}}$ with $n=20$ (left) $n=40$ (center) and $n=80$ (right).}
% \label{funf123n}
% \end{figure}


% To better appreciate the quality of the approximation operators that we have proposed, we reconstruct the functions $f_1$–$f_3$ using the histopolation-regression operator based on Padua triangles relative to the regular triangulation $\T_{20,20}^{\mathrm{reg}}$. These reconstructions are displayed in Fig.~\ref{rec}.


% \begin{figure}[ht]
%     \centering
%         \includegraphics[scale=0.32]{f1rec.eps}
%     \includegraphics[scale=0.32]{f2rec.eps}
%     \includegraphics[scale=0.32]{f3rec.eps}
%     \caption{Reconstruction of the functions $f_1$ (left), $f_2$ (center) and $f_3$ (right) produced by $\Pi^{\mathrm{Pad}}_m$ over the domain $\Omega=[-1,1]^2$ using the regular triangulation $\T_{20,20}^{\mathrm{reg}}$.}
% \label{rec}
% \end{figure}


% In all the experiments presented so far, we have considered domains partitioned by triangulations whose triangles are not excessively distorted, i.e., they do not have excessively \textit{small} angles. However, when this condition is not satisfied, the strategies discussed here may yield not good results. To prove this, we examine the behavior of $\hat{e}_d^{\mathrm{Fek}}\left(f_1\right)$ when the domain $\Omega=[-1,1]^2$ is partitioned using general triangulations consisting of $N$ triangles, with $N$ ranging from $82$ to $402$, see Fig~\ref{trianacut}. The result of this experiment is shown in Fig~\ref{errorsss}.

% \begin{figure}[ht]
%     \centering
%     \includegraphics[scale=0.32]{122triangles.eps}
%     \includegraphics[scale=0.32]{202triangles.eps}
%     \includegraphics[scale=0.32]{362triangles.eps}
%     \caption{Plot of a  triangulation consisting of $N=122$ (left) $N=202$ (center) and $N=362$ (right) triangles over the domain $\Omega = [-1,1]^2$. The red triangles represent the corresponding Padua triangles.}
% \label{trianacut}
% \end{figure}




% \begin{figure}[ht]
%     \centering    \includegraphics[scale=0.49]{errorrand.eps}
%     \caption{Semilog plot of the maximum approximation error produced by $\hat{\Pi}^{\mathrm{Fek}}_m$ over the domain $\Omega=[-1,1]^2$, using general triangulations with the number of triangles ranging from $82$ to $402$}
% \label{errorsss}
% \end{figure}

% From this plot, we observe that the error does not decrease as the number of triangles in the triangulation increases. Therefore, we conclude that convergence is not achieved. Instead, we observe that the error decreases up to a sufficiently large number of triangles, after which the error begins to grow dramatically.


\section*{Declarations}

\vspace{-0.2cm}

\begin{itemize}
    \item Funding: This research was supported by GNCS-INdAM 2024 projects. The first author is funded by IN$\delta$AM and supported by the University of Padova. The third and the fourth authors are funded by the European Union – NextGenerationEU under the National Recovery and Resilience Plan (NRRP), Mission 4 Component 2 Investment 1.1 - Call PRIN 2022 No. 104 of February 2, 2022 of the Italian Ministry of University and Research; Project 2022FHCNY3 (subject area: PE - Physical Sciences and Engineering) \lq\lq Computational mEthods for Medical Imaging (CEMI)\rq\rq. 
    \item Authors' Contributions: All authors contributed equally to this article, and therefore the order of authorship is alphabetical.
    \item Conflicts of Interest: The authors declare no conflicts of interest.
    \item Ethics Approval: Not applicable.
    \item Data Availability: No data were used in this study.
\end{itemize}

\vspace{-0.6cm}



%\section*{Acknowledgments}
% This research has been conducted as part of RITA \textquotedblleft Research
% ITalian network on Approximation'' and as part of the UMI group \enquote{Teoria dell'Approssimazione
% e Applicazioni}. The research was supported by GNCS-INdAM 2025 projects. The first author is funded by IN$\delta$AM and supported by the University of Padova. The third and the fourth author are funded by the European Union – NextGenerationEU under the National Recovery and Resilience Plan (NRRP), Mission 4 Component 2 Investment 1.1 - Call PRIN 2022 No. 104 of February 2, 2022 of Italian Ministry of University and Research; Project 2022FHCNY3 (subject area: PE - Physical Sciences and Engineering) \enquote{Computational mEthods for Medical Imaging (CEMI)}.


\bibliographystyle{spmpsci}
\bibliography{paper}





\end{document}
